% Modernised reading — fonds Alexandre Grothendieck, shelfmark 162-5
% (pages 1-68, the whole folder).
%
% This is an INTERPRETATION, not a transcription. It re-reads the pages in
% today's notation and vocabulary, states the hypotheses the pages leave
% implicit, and drops the critical apparatus so that the mathematics reads
% straight through; everything the apparatus carried moves into footnotes. It is
% held to being mathematically correct as it stands: where the pages are loose
% or mistaken, the true statement is what appears here and a footnote says what
% the page has. In any dispute of substance the transcription governs: whoever
% cites Grothendieck must cite batch-01.fr.tex (pp. 1-20), batch-02.fr.tex
% (pp. 21-40) or batch-03.fr.tex (pp. 41-60). Pages 61-68 have no transcription,
% by design: they are pp. 5-12 of Quillen's spirit-duplicated English typescript
% « On a conjecture of Adams », not his, carrying no annotation.
%
% Pass: Opus 5.5 (claude-opus-5-5), 2026-09-22 — first pass, unchecked against the pages by a human.
% The folder was read whole, all three transcribed batches together; this is
% one document for the shelfmark. It was made on Opus 5.5 and has not been
% measured against the reference reading of folder 115 (made on Opus 5).
%
% Witnesses. The typescript « Tapis de Quillen » (notes du 10.9.68, 16 typed
% pages) is here twice. Copy B (pp. 32-41, 43-48) is the typed top copy
% annotated in blue ink; its transcription (batches 2-3) is the text followed
% here for the typescript AND for its ink layer. Copy A (pp. 2-17) is a black
% photocopy of copy B taken after the ink, which then received notes of its
% own. From batch 1 this reading takes only: the cover (p. 1), the date « notes
% du 10.9.68 », and copy A's own notes — « (compactes ?) » (p. 12), « NB H*(X)
% (= H*(X, F_2)) » (p. 14), « i.e. l'enveloppe "karoubienne" » (p. 14),
% « quasi-abélienne = karoubienne + additive » (p. 15), « cf exposé Karoubi à
% Bourbaki ! » (p. 17). Batch 1 was being revised when this was written; its
% readings of the ink layer shared with copy B are NOT used. Statements that
% rest on copy A: the compactness footnote (part III), the F_2-coefficients
% footnote, the two Karoubi-envelope footnotes, and the dating inference from
% the Karoubi pointer (section on the photocopy). Folder 111 is a third copy of
% the typescript; it is named, not used.
%
% Spine. pp. 21-30 (6.9.68): formal categories over X_0, abstract De Rham
% complexes, the Lie algebroid g = Hom(Omega, O), Quillen's theorem in char. 0,
% coefficients, De Rham vs Hochschild. pp. 32-48 (10.9.68): the typescript,
% parts I (categories and simplicial sets), II (n-categories, Gr-categories,
% Dold-Kan, K_n), III (cobordism « motivique »). The typescript takes up NOTHING
% of the manuscript, in particular not its section 11 (De Rham-Hochschild).
% Order of the body: manuscript, typescript (copy B), photocopy (copy A),
% printed pieces by Quillen (pp. 19, 50-68).
%
% Notation fixed for the whole document. Manuscript: O = O_{X_0}; C a formal
% category, A its pro-ring, I the augmentation ideal, Omega = I/I^2; delta for
% the differential of any De Rham complex (the page's d = d_C when it comes from
% C); g = Hom_O(Omega, O), theta its anchor, i_X contraction; pairing
% <a^b, X^Y> = a(X)b(Y) - a(Y)b(X), determinants on Lambda^3; nabla for the
% connection M -> Omega (x) M (the page's delta^1_M on p. 27, delta^0_M on
% p. 28). Typescript: N(C) for the nerve (the page's S(C)); Ho(Cat), Ho(sSet)
% for (Cat)', (Ssimpl)'; C^ for the presheaf topos (the page's C~); « catégorie
% double » for the page's « bi-catégorie »; V for the category of manifolds of
% part III (the page's C); B_d(X) graded by relative dimension; H^k(X) =
% H^k(X; F_2), matched with B_{-k}; Hom_{B^H} for the page's H(X,Y).
%
% Substantive departures, with pages.
%  p. 22-23  Remarque 1 claims alpha), beta), gamma) suffice for delta^2 = 0.
%            False: delta^2 on Omega is a further O-linear condition,
%            equivalent (Omega loc. free f.t.) to Jacobi. Counterexample: X_0 a
%            point, any antisymmetric non-Jacobi bracket. beta) corrected
%            (delta_0 f ^ omega).
%  p. 24-25  « pourquoi est-ce nul ?? »: answered (ours) by the identity
%            <omega, J(X,Y,Z)> = <delta delta omega, X^Y^Z>
%              - sum_cyc <delta delta <omega,Z>, X^Y>,
%            checked by hand and numerically; the page's route through
%            T(delta omega) is not followed.
%  p. 25-26  §5 correspondence and §6 Quillen's theorem: variance fixed
%            (covariant in A, contravariant in C; C -> g^C covariant).
%  p. 28     sign (-1)^{deg omega} supplied in the extended Leibniz rule.
%  p. 29     K(X,Y)(u) = -<K'(u), X^Y> (the page: K'(u)(X,X), « Attention
%            signe ! »).
%  p. 30     descent data <-> INTEGRABLE connections (the page: connexions).
%  p. 32     « système de coefficients » = functors C° -> Ens inverting every
%            arrow (the page: « transformant isomorphismes en isomorphismes »,
%            which every functor does); lim^(i)_C F read as lim^i over C° of
%            u*F'.
%  p. 34     coproduct of the cochains corrected to a product (the ink asks it).
%  p. 35     C, C' read X, X'; the question « se borner aux systèmes locaux
%            commutatifs » answered negatively (acyclic groups, ours).
%  p. 36     unit functor C_0 -> C_1 supplied; Segal condition supplied for the
%            page's « objets qui … ».
%  p. 38     Hom_grab identified with tau_{<=0} RHom (truncation supplied).
%  p. 39     the ink « pl. fidèle (?) » is corrected: the loop functor is not
%            even faithful (cup-product example, ours).
%  p. 39     « tuage » analogy for Aut(1): footnoted as rather a looping.
%  p. 43     compact (closed) manifolds supplied.
%  p. 44     F^• : V° -> B (the page: C -> B); composition in B written out.
%  p. 45     base change written alpha_{s•} alpha_s^• = Gamma^• i_{s•} (the
%            page's order alpha_s^* alpha_{s*} does not type-check).
%  p. 45     H (x) Lambda = B holds as Lambda-modules for any homogeneous lift,
%            by Conner-Floyd freeness and graded Nakayama.
%  p. 45-46  the page's « quasi-abélienne » = pseudo-abelian (Karoubian
%            additive), not Schneiders' quasi-abelian.
%  p. 47-48  « Thom's theorem » (x in S(X) iff A^+ x = 0), read with the usual
%            action on cohomology, fails: RP^1 in RP^2. Reported under reserve;
%            the page itself doubts the necessity.
%
% Announced and never established on the pages: the divided-power variant
% (pp. 26, 30); whether a formal category is a groupoid (p. 21); whether order 2
% suffices in the descent condition (p. 27); the lemma of p. 23; the
% generalisation to « localement homotopiquement trivial » topoi (pp. 34-35);
% n-categories as n-simplicial sets (p. 37); the Hom_grab identification
% (p. 38) and the 2-Grab-category of homomorphisms (p. 40); abelian groups in
% n-Cat <-> complexes of length n (p. 40); K_n via universal Gr-n-categories
% (p. 40); the universality of S (p. 47); the structure of G (p. 48).
\documentclass[11pt,a4paper]{article}
% Le préambule commun à toutes les transcriptions du fonds Grothendieck.
%
% Il tient en une idée : **l'incertitude se compose, elle ne se résout pas.**
% Une transcription qui lisse un mot douteux a détruit la seule chose pour
% laquelle on la faisait — un lecteur doit pouvoir savoir, à chaque endroit,
% ce qui était sur le feuillet et ce qui a été ajouté. Les six macros qui
% suivent sont donc l'appareil critique, et non de la décoration.
%
% Les mêmes commandes sont reconnues par `scripts/render.mjs`, qui produit la
% vue de lecture du site. Une macro ajoutée ici doit l'être aussi là-bas,
% faute de quoi le rendu échouera bruyamment — ce qui est voulu : un rendu
% silencieusement faux est pire qu'un rendu absent.

\ProvidesPackage{grothendieck}[2026/08/08 Transcriptions du fonds Grothendieck]

\usepackage{amsmath,amssymb,amsthm}
\usepackage{mathrsfs}
\usepackage{stmaryrd}
% Les diagrammes commutatifs sont partout dans ce fonds : c'est la langue
% même de la géométrie algébrique de Grothendieck, et une transcription qui
% les rendrait en prose ne transcrirait rien.
\usepackage{tikz-cd}
% Français par défaut — c'est la langue des feuillets — mais l'anglais est
% chargé aussi : la traduction et le résumé sont des documents anglais, et la
% typographie française (espaces avant les deux-points) y serait une faute.
% Ces documents basculent par \selectlanguage{english} en tête de corps.
\usepackage[english,french]{babel}
\usepackage{fontspec}
\usepackage{geometry}
\geometry{a4paper,margin=25mm,marginparwidth=28mm}
\usepackage{marginnote}
\usepackage{xcolor}
% ulem, not soul. `soul`'s \st and \ul are built on a character-by-character
% scan that XeLaTeX's font machinery breaks: the document fails with an
% undefined control sequence rather than degrading. ulem does the same two jobs
% and is Unicode-engine safe. `normalem` stops it hijacking \emph, which the
% transcriptions use for Grothendieck's own emphasis.
\usepackage[normalem]{ulem}
\usepackage{hyperref}
\hypersetup{colorlinks=true,linkcolor=black,urlcolor=black,pdfborder={0 0 0}}

% --- Métadonnées du lot -----------------------------------------------------
% Reprises telles quelles par le script de rendu, qui les affiche en tête de
% la vue de lecture. Elles fixent ce que le fichier prétend couvrir : sans
% elles, une transcription flotte sans point d'attache dans un fonds de seize
% mille feuillets.

\newcommand{\folder}[1]{\gdef\grothFolder{#1}}
\newcommand{\batch}[1]{\gdef\grothBatch{#1}}
\newcommand{\leaves}[2]{\gdef\grothFirst{#1}\gdef\grothLast{#2}}
\newcommand{\foldertitle}[1]{\gdef\grothFoldertitle{#1}}
\gdef\grothFolder{?}\gdef\grothBatch{?}\gdef\grothFirst{?}\gdef\grothLast{?}\gdef\grothFoldertitle{}

% --- Le résumé --------------------------------------------------------------
%
% Il ouvre la lecture modernisée, sous le titre « Résumé », et il est composé
% en petit. Ce n'est pas un ornement : c'est le seul passage du document qui
% ne soit pas des mathématiques, et un lecteur doit pouvoir le reconnaître
% avant de l'avoir lu — puis le sauter s'il connaît déjà le sujet, ou s'y
% arrêter s'il ne le connaît pas. Le filet qui le suit marque la reprise du
% corps.

\newenvironment{resume}
  {\small\setlength{\parskip}{0.45em}}
  {\par\medskip\hrule\medskip}

% --- Mots-clés ---------------------------------------------------------------
%
% En anglais, en clôture du résumé de la lecture modernisée : le vocabulaire
% moderne sous lequel chercher ce que les feuillets construisent. Le manifeste
% du site (`npm run manifest`) les extrait pour étiqueter la cote entière —
% cette ligne est la source unique des tags, il n'y a pas de fichier à côté.
\newcommand{\keywords}[1]{\par\smallskip\noindent
  {\footnotesize\textsc{Keywords} --- \textit{#1}}\par}

% --- L'appareil critique ----------------------------------------------------

% Le numéro de feuillet, en marge. C'est l'ancre qui permet de régler un
% désaccord sur une formule en regardant *un* feuillet plutôt qu'une chemise
% de six cents. Le site s'en sert aussi pour tourner les pages du fac-similé
% quand on fait défiler la transcription.
\newcommand{\leaf}[1]{%
  \marginnote{\footnotesize\color{gray!70}\textsf{#1}}%
  \label{leaf:#1}\ignorespaces}

% Illisible : on ne devine pas. Un mot inventé qui se lit comme les autres est
% le pire résultat possible.
\newcommand{\ill}{\textcolor{red!60!black}{[\,\ldots\,]}}

% Lecture douteuse : le mot est donné, et signalé comme incertain.
\newcommand{\uncertain}[1]{\uline{#1}}

% Ajout de l'éditeur — une lettre manquante, un mot sous-entendu.
\newcommand{\add}[1]{\textcolor{blue!50!black}{[#1]}}

% Biffé par Grothendieck lui-même. Ce qu'il a rayé fait partie du document :
% une rature dit souvent où la pensée a changé de direction.
\newcommand{\struck}[1]{\sout{#1}}

% Note du transcripteur, sur l'état matériel du feuillet.
\newcommand{\note}[1]{\par\smallskip{\small\color{gray!60!black}\textit{[#1]}}\par\smallskip}

% Note marginale de Grothendieck, distincte du corps du texte.
\newcommand{\marginal}[1]{\par\smallskip\noindent\hspace{1em}%
  {\small\color{gray!60!black}#1}\par\smallskip}

% --- Titre ------------------------------------------------------------------

\AtBeginDocument{%
  \begin{center}
    {\large\bfseries Fonds Alexandre Grothendieck --- cote n\textsuperscript{o}~\grothFolder}\\[2pt]
    {\itshape\grothFoldertitle}\\[4pt]
    {\small feuillets \grothFirst--\grothLast\ (lot \grothBatch)}\\[2pt]
    {\footnotesize\color{gray!60!black}Transcription établie d'après le fac-similé numérisé
     par l'Université de Montpellier.\\
     Les lectures incertaines sont soulignées, les illisibles notées [\,\ldots\,],
     les ajouts entre crochets.}
  \end{center}
  \vspace{1em}}

\endinput


\folder{162-5}
\pages{1}{68}
\dating{1968-[à partir de 1970] — le groupe « Dossiers rassemblés par Grothendieck sur différentes thématiques » (161-1 à 162-6) est daté [après 1961-vers 1977]}
\watermark{Édition de démonstration\\interprétation personnelle de l'œuvre}
\foldertitle{Tapis de Quillen : tapuscrit et copies de tapuscrits (1968, s.d.), notes manuscrites (1968), tiré à part (1968) — lecture modernisée du dossier entier}

\begin{document}

\section*{Résumé}

\begin{resume}
Au début de septembre 1968, Grothendieck a pris des notes sur ce que Daniel
Quillen, alors âgé de vingt-huit ans, lui exposait. Ce dossier les réunit.
« Tapis », dans sa langue de ces années, désigne un ensemble de résultats et la
manière de s'en servir : les mêmes pages parlent du « tapis de Dold » et du
« tapis de Roos ». Il y a deux textes, écrits à quatre jours d'intervalle et sur
des sujets qui ne se recoupent pas. Le premier, du 6 septembre, fait dix pages
manuscrites. Le second, du 10, est un tapuscrit de seize pages, présent ici en
deux exemplaires annotés. Autour, des pièces de Quillen lui-même : une liste de
ses publications, un tiré à part, une prépublication en anglais.

Les notes du 6 septembre portent sur les symétries infinitésimales. Un groupe
de Lie a une algèbre de Lie, qui en est comme l'ombre au premier ordre : elle
retient les directions dans lesquelles on peut partir de l'identité, et la façon
dont deux de ces mouvements ne commutent pas. En caractéristique nulle, cette
ombre suffit à reconstruire le groupe au voisinage de l'identité. Le théorème
que les notes attribuent à Quillen dit la même chose pour des objets plus
souples, des « catégories formelles ». Au lieu d'un seul point, on a tous les
points d'un espace, et entre deux points voisins des flèches infiniment
petites. L'ombre au premier ordre d'un tel objet est ce qu'on appelle
aujourd'hui un algébroïde de Lie. Grothendieck le code par une algèbre
différentielle graduée, un « complexe de De Rham » généralisé. Suivent les
coefficients : un module muni d'une action de la catégorie formelle est la même
chose, en caractéristique nulle, qu'un module muni d'une connexion plate. La dernière page compare deux
cohomologies de ce module, et c'est là que les notes s'arrêtent. Une question
y est laissée ouverte, « pourquoi est-ce nul ?? », sur l'identité de Jacobi.
Elle a une réponse simple, que nous donnons ; elle oblige à corriger une remarque
antérieure des mêmes pages.

Le tapuscrit du 10 septembre a trois parties. La première part de ce qu'une
catégorie est un espace : ses objets sont des points, ses flèches des segments,
ses triangles commutatifs des faces. Selon Quillen, les catégories et les
ensembles simpliciaux ont la même théorie de l'homotopie, et toute catégorie a
le même type d'homotopie que son opposée. Grothendieck y voit l'occasion de
définir un type d'homotopie par la catégorie dérivée de ses systèmes locaux.
La deuxième partie traite des catégories supérieures munies d'une loi de groupe.
Une catégorie de Picard, c'est-à-dire un groupe « à isomorphismes près », est
à peu de chose près un complexe de groupes abéliens de longueur un. Il en va de
même en longueur deux, et le texte espère qu'il en ira ainsi en toute longueur.
Il espère aussi définir les groupes $K_n$ d'une catégorie munie d'un produit
tensoriel comme les groupes d'homotopie de sa « complétion en groupe »
universelle. C'est ce que Quillen fera dans les années suivantes, par d'autres
chemins. La troisième partie regarde le cobordisme d'un œil de géomètre
algébriste. Deux variétés sont cobordantes quand, ensemble, elles bordent une
troisième. En comptant les variétés au-dessus d'un espace à cobordisme près, on
obtient une catégorie de correspondances qui ressemble aux motifs. Le groupe
des symétries des opérations de Steenrod y joue le rôle d'un « groupe de Galois
motivique ». Contrairement au cas familier, ce groupe est unipotent, et il est
défini sur le corps à deux éléments.

Ces pages montrent aussi une manière de travailler. Le tapuscrit rapporte
(« Quillen prouve ») et se pose des questions (« Demander à Quillen… »). Dans ses
marges, une encre différente répond à certaines d'entre elles, par « oui » ou
par « non, … ». On y voit une conversation mathématique en train de se
faire : ce qu'on a compris, ce qu'on croit, ce qu'on ira demander.

Pour aller plus loin, les noms à chercher sont les suivants : algébroïdes de
Lie et groupoïdes formels ; la structure de modèles de Thomason sur les
catégories ; les 2-groupes et les champs de Picard ; la $K$-théorie supérieure
de Quillen. Il faut surtout lire la note de Quillen de 1969 sur les lois de
groupes formels du cobordisme, qui, un an après ces notes, donnera à la
troisième partie sa structure véritable. On la retrouve aujourd'hui dans le
cobordisme algébrique de Levine et Morel.

\keywords{Lie algebroid, formal groupoid, Chevalley–Eilenberg complex, integrable connection, van Est isomorphism, nerve of a category, category of simplices, homotopy theory of categories, double category, Segal condition, Dold–Kan correspondence, 2-group, Picard category, higher algebraic K-theory, group completion, unoriented cobordism, correspondence category, Karoubi envelope, Chow motive, Steenrod algebra, motivic Galois group, Pontryagin–Thom construction}
\end{resume}

\pagerange{1}{68}
\section*{Le fil du dossier, et les conventions}

\emph{Les pièces}, dans l'ordre des archivistes.

\begin{enumerate}
  \item \emph{Page 1}, la chemise : « Le Tapis de Quillen / notes d'Alexandre
  (1968) », au crayon, d'une main qui n'est probablement pas la sienne.
  \item \emph{Pages 2 à 17} : une photocopie noire du tapuscrit « Tapis de
  Quillen » (ses pages 1 à 16), datée « notes du 10.9.68 ». Nous l'appelons la
  \emph{photocopie}.
  \item \emph{Page 19} : une liste dactylographiée, « Publications of Daniel
  Quillen », sans annotation. Les pages 18 et 20 sont blanches.
  \item \emph{Pages 21 à 30} : dix feuillets manuscrits à l'encre bleue,
  intitulés « tapis Quillen (Notes du 6.9.68) » et numérotés par paragraphes
  de 1 à 11. La page 31 est blanche.
  \item \emph{Pages 32 à 41 et 43 à 48} : l'original dactylographié du même
  tapuscrit, annoté à l'encre bleue. Nous l'appelons l'\emph{exemplaire annoté}.
  La page 42 est blanche.
  \item \emph{Pages 50 à 55} : le tiré à part de D. Quillen, « On the
  Representation of Hermitian Forms as Sums of Squares », Inventiones math. 5
  (1968). Les pages 49 et 56 sont blanches.
  \item \emph{Pages 57 à 68} : un tapuscrit anglais de Quillen reproduit à
  l'alcool, « On a conjecture of Adams », ses pages 1 à 12. Seule la page 58 y
  porte une intervention étrangère à la reproduction.
\end{enumerate}

\emph{Les deux exemplaires du tapuscrit.} La photocopie a été tirée sur
l'exemplaire annoté \emph{après} que l'encre y eut été portée. Elle en reproduit
donc la couche manuscrite, puis elle a reçu quelques notes qui lui sont propres.
Le texte suivi ici, pour le tapuscrit comme pour ses annotations, est celui de
l'exemplaire annoté. Les notes propres à la photocopie sont signalées en note au
point où elles portent, et rassemblées dans une section à part. Le dossier 111
du fonds contient une troisième copie du même tapuscrit.

\emph{Deux textes sans intersection.} Les notes du 6 septembre et le tapuscrit
du 10 portent le même titre, mais le tapuscrit ne reprend rien des notes. Il n'y
est question ni de catégories formelles, ni de complexes de De Rham, ni, en
particulier, de la comparaison entre cohomologie de De Rham et cohomologie de
Hochschild par laquelle les notes se terminent (page 30). Ce sont deux
conversations, ou deux moments d'une conversation, sur des sujets distincts. On
les lit ici dans l'ordre des dates.

\emph{Qui écrit.} Le tapuscrit parle à la première personne : « ma lettre à
Verdier », « “ma” construction chérie de la “catégorie de Picard” », « mon
interprétation chérie », « le rôle universel que j'ai envisagé plus haut ». Il
rapporte ce que « Quillen prouve » et note ce qu'il faudra « demander à
Quillen ». Nous n'attribuons pas l'encre qui annote les deux exemplaires ; la
transcription de l'exemplaire annoté ne le fait pas non plus. Les notes du 6
septembre sont de sa main selon la transcription.

\emph{Notations}, valables pour tout le document ; chaque écart avec la page est
signalé en note.

\begin{itemize}
  \item Notes du 6 septembre. $X_0$ est un schéma et $\mathcal{O} =
  \mathcal{O}_{X_0}$. On note $\mathcal{C}$ une catégorie formelle sur $X_0$, $A$
  son pro-anneau, $I$ l'idéal d'augmentation, $\Omega = I/I^{2}$ le conormal et
  $\Lambda^{\bullet}\Omega$ l'algèbre extérieure sur $\mathcal{O}$. La
  différentielle d'un complexe de De Rham est notée $\delta$ ; la page écrit $d$
  quand elle vient de $\mathcal{C}$. On pose $\mathfrak{g} =
  \mathcal{H}\mathrm{om}_{\mathcal{O}}(\Omega, \mathcal{O})$, on note $\theta$ son
  ancre et $i_X$ la contraction par $X$. L'accouplement est $\langle a \wedge b, X
  \wedge Y\rangle = a(X)b(Y) - a(Y)b(X)$, par déterminants en degré 3. Pour un
  module $M$, une connexion est notée $\nabla : M \to \Omega \otimes M$.
  \item Tapuscrit. $N(C)$ désigne le nerf (la page : $S(C)$), et $C^{\circ}$ la
  catégorie opposée. $\mathrm{Ho}(\mathrm{Cat})$ et $\mathrm{Ho}(\mathrm{sSet})$
  désignent les catégories homotopiques (la page : $(\mathrm{Cat})'$,
  $(\mathrm{Ssimpl})'$), et $\widehat{C}$ le topos des préfaisceaux sur $C$ (la
  page : $\widetilde{C}$). Dans la troisième partie, $\mathcal{V}$ est la
  catégorie des variétés (la page : $C$), et $H^{\bullet}(X) = H^{*}(X;
  \mathbb{F}_2)$.
\end{itemize}

\emph{Ce que le dossier annonce sans l'établir.} La variante à puissances
divisées du théorème de Quillen est annoncée deux fois (pages 26 et 30) et
jamais écrite. Plusieurs questions restent ouvertes : si une catégorie formelle
est automatiquement un groupoïde (page 21), et si l'ordre 2 suffit dans la
condition de descente (page 27). Le tapuscrit laisse ouverte la généralisation
aux topos « localement homotopiquement triviaux » (pages 34 et 35). Il pose sans
les établir la description des $n$-catégories par des ensembles
$n$-simpliciaux (page 37) et celle des homomorphismes entre catégories de
Picard (pages 38 et 40). Il ne fait qu'espérer la correspondance entre groupes
abéliens de $(n\text{-Cat})$ et complexes de longueur $n$ (page 40), et la
définition des $K_n$ (page 40). Il suppose sans la démontrer l'universalité de
la catégorie $\mathbf{S}$ (page 47), et ne dit pas la structure du groupe $G$
(page 48).

\pagerange{21}{30}
\section*{Les notes du 6 septembre 1968 : catégories formelles et complexes de De Rham (pages 21 à 30)}

\pagerange{21}{22}
\subsection*{1. Catégories formelles (pages 21 et 22)}

Soit $X_0$ un schéma. Une \emph{catégorie formelle} $\mathcal{C}$ sur $X_0$ est
un objet catégorie dans les schémas formels, c'est-à-dire dans un certain type
d'ind-schémas. Son schéma d'objets est $X_0$ lui-même, et son schéma de flèches
$X_1$ est un voisinage infinitésimal de $X_0$, plongé par les flèches
identiques.\footnote{En marge gauche de la page 21, le schéma vertical d'un objet
semi-simplicial tronqué $X_2 \to X_1 \rightrightarrows X_0$ : trois flèches de
faces descendent de $X_2$ à $X_1$, deux de $X_1$ à $X_0$, avec les
dégénérescences en sens inverse. La page regarde $X_0$ « comme schéma formel,
i.e.\ ind-schéma particulier ».} Algébriquement, c'est la donnée des objets
suivants :
\begin{itemize}
  \item un pro-anneau $A = (A_\alpha)$ sur $X_0$, qualifié d'« admissible » ;
  \item une augmentation $\varepsilon : A \to \mathcal{O}$, dont l'idéal $I$
  est un idéal de définition ;
  \item deux structures de $\mathcal{O}$-algèbre, à gauche et à droite, qui
  correspondent à la source et au but ;
  \item une comultiplication $\Delta : A \to A \otimes_{\mathcal{O}} A$, le
  produit tensoriel étant pris entre la structure droite du premier facteur et
  la structure gauche du second.
\end{itemize}
On demande que $\Delta$ soit coassociative et que l'augmentation soit une
counité à gauche et à droite. On dirait aujourd'hui un groupoïde formel, ou
une algèbre de Hopf complète au-dessus de $\mathcal{O}$ (un « algébroïde de
Hopf »), mais sans antipode. La page demande si l'antipode existe
nécessairement, c'est-à-dire si une catégorie formelle est automatiquement un
groupoïde formel.\footnote{Question laissée ouverte par la page, entre
crochets ; les mots « de passage à l'inverse » y sont d'une lecture incertaine.
Pour une loi de groupe formel sur un anneau quelconque, l'inverse existe
toujours, ce qui rend plausible une réponse positive au moins dans le cas lisse.
Nous ne tranchons pas.}

On pose $\Omega_{\mathcal{C}} = I/I^{2}$, le faisceau conormal de $X_0$ dans
$X_1$. C'est a priori un pro-faisceau quasi-cohérent sur $X_0$, et un module
quasi-cohérent ordinaire lorsque $A$ est $I$-adique. Sur
$\Lambda^{\bullet}\Omega_{\mathcal{C}}$ on définit une différentielle d'anneau
gradué $\delta = d_{\mathcal{C}}$, dont la composante de degré 0 est
\[
\delta f = t^{*}f - s^{*}f \mod I^{2},
\]
où $s, t : X_1 \to X_0$ sont la source et le but (la page écrit
$\mathrm{pr}_1$, $\mathrm{pr}_2$). La construction en degrés supérieurs est
renvoyée « plus bas » ; elle est ce que le paragraphe suivant rend abstrait.

\emph{Opérateurs différentiels.} Regardons $A_\alpha$ comme
$\mathcal{O}$-module par la structure gauche. L'application $f \mapsto t^{*}f$
de $\mathcal{O}$ dans $A_\alpha$ est alors un opérateur différentiel, d'ordre
$n$ si $I_\alpha^{n+1} = 0$. Les sections de
$\mathcal{D}^{\mathcal{C}}_\alpha =
\mathcal{H}\mathrm{om}_{\mathcal{O}}(A_\alpha, \mathcal{O})$ opèrent donc sur
$\mathcal{O}$ par opérateurs différentiels. La réunion
$\mathcal{D}^{\mathcal{C}} = \varinjlim_\alpha \mathcal{D}^{\mathcal{C}}_\alpha$
est un faisceau d'anneaux associatifs et unitaires, dont le produit est dual de
$\Delta$. Elle s'envoie dans le faisceau des opérateurs différentiels de
Grothendieck $\mathcal{D}_{X_0/\mathbb{Z}}$. Le cas universel est celui où
$\mathcal{C}$ est le complété formel de la diagonale de $X_0 \times_{\mathbb{Z}}
X_0$. Les $A_\alpha$ y sont alors les faisceaux de parties principales, et
$\mathcal{D}^{\mathcal{C}}$ est $\mathcal{D}_{X_0/\mathbb{Z}}$
lui-même.\footnote{La page le dit dans une parenthèse coupée par le changement
de feuillet, avec un mot illisible : « au fond ici ce serait la formalisation de
la catégorie … définie par $X_0/\mathrm{Spec}\,\mathbb{Z}$ ». L'identification
aux parties principales et aux opérateurs différentiels de EGA IV, § 16, est
notre explicitation.} En particulier, le module
\[
\mathfrak{g}^{\mathcal{C}} =
\mathcal{H}\mathrm{om}_{\mathcal{O}}(\Omega_{\mathcal{C}}, \mathcal{O})
= \varinjlim_\alpha \mathcal{H}\mathrm{om}_{\mathcal{O}}(\Omega_{\mathcal{C}_\alpha}, \mathcal{O})
\]
opère sur $\mathcal{O}$ par dérivations. C'est le module des « champs de
vecteurs » de $\mathcal{C}$. On le retrouvera au paragraphe 3 à partir du seul
complexe de De Rham.

\pagerange{22}{23}
\subsection*{2. Complexes de De Rham abstraits (pages 22 et 23)}

Soit $\Omega$ un $\mathcal{O}$-module quasi-cohérent. Un \emph{complexe de De Rham}
sur $\Omega$ est une différentielle d'anneau gradué $\delta$ sur
$\Lambda^{\bullet}\Omega$. C'est une antidérivation de degré $+1$,
$\delta(ab) = \delta a \cdot b + (-1)^{|a|} a \cdot \delta b$, de carré
nul.\footnote{« abstrait » est d'une lecture incertaine. La page ajoute entre
crochets qu'il faudrait considérer aussi les espaces annelés, ainsi que « tout
complexe obtenu par les sorites de 1 » ; c'est le cas du complexe de la
catégorie formelle.} Autrement dit, c'est une structure d'algèbre différentielle
graduée commutative sur $\Lambda^{\bullet}\Omega$, compatible à la graduation.

\emph{Remarque 1, rectifiée.} Une telle $\delta$ est déterminée par ses deux
premières composantes. La première est
$\delta_0 : \mathcal{O} \to \Omega$, une dérivation :
\[
\alpha)\quad \delta_0(fg) = \delta_0(f)\,g + f\,\delta_0(g).
\]
Il revient au même de se donner un morphisme $u : \Omega^{1}_{X_0/\mathbb{Z}} \to
\Omega$, avec $\delta_0 = u \circ d$. La seconde est $\delta_1 : \Omega \to
\Lambda^{2}\Omega$, soumise aux conditions
\[
\beta)\quad \delta_1(f\omega) = \delta_0(f) \wedge \omega + f\,\delta_1(\omega),
\qquad
\gamma)\quad \delta_1\delta_0 = 0 .
\]
Inversement, soient $\delta_0$, $\delta_1$ vérifiant $\alpha)$ et $\beta)$.
Ils se prolongent de façon unique en une antidérivation $\delta$ de
$\Lambda^{\bullet}\Omega$.\footnote{La page écrit, au premier terme de
$\beta)$, $\delta_0(f)\,\delta_1(\omega)$ ; la règle de Leibniz donne
$\delta_0(f) \wedge \omega$, que nous rétablissons.} Le carré
$\delta^{2}$ est alors une dérivation de degré 2, donc nul si et seulement s'il
l'est sur $\mathcal{O}$ et sur $\Omega$. Sur $\mathcal{O}$, c'est la condition
$\gamma)$. Sur $\Omega$, la condition $\gamma)$ rend
$\delta^{2}|_{\Omega} : \Omega \to \Lambda^{3}\Omega$ $\mathcal{O}$-linéaire, car
$\delta^{2}(f\omega) = \delta^{2}f \wedge \omega + f\,\delta^{2}\omega$. Mais elle
ne l'annule pas : c'est une condition de plus,
\[
\gamma')\quad \delta_2\,\delta_1 = 0 ,
\]
où $\delta_2$ est la composante de degré 2 de l'antidérivation prolongée. On
verra au paragraphe suivant qu'elle équivaut à l'identité de Jacobi, lorsque
$\Omega$ est localement libre de type fini.\footnote{La page affirme que
l'antidérivation prolongée « est de carré nul grâce à la condition
$\gamma)$ », et note en marge : « dégager un lemme circonstancié ». C'est faux
tel quel. Soit $X_0$ un point, $\Omega = V^{*}$ le dual d'un espace vectoriel
de dimension finie, $\delta_0 = 0$ et $\delta_1$ la transposée d'une
application bilinéaire alternée quelconque $V \times V \to V$. Les conditions
$\alpha)$, $\beta)$ et $\gamma)$ sont trivialement satisfaites, et $\delta^{2} =
0$ sur $\Omega$ équivaut exactement à l'identité de Jacobi pour ce crochet. C'est
elle que la page cherchera en vain à démontrer aux pages 24 et 25. La définition
elle-même, au début du paragraphe, demande bien une différentielle,
donc $\delta^{2} = 0$ : seule la Remarque est en défaut.}

\pagerange{23}{25}
\subsection*{3. L'algébroïde de Lie d'un complexe de De Rham (pages 23 à 25)}

Soit $(\Lambda^{\bullet}\Omega, \delta)$ un complexe de De Rham, et
$\mathfrak{g} = \mathcal{H}\mathrm{om}_{\mathcal{O}}(\Omega, \mathcal{O})$.
C'est un faisceau de $\mathcal{O}$-modules.\footnote{La page le dit
« ind-module quasi-cohérent, qui s'identifie à un Module quasi-cohérent bien
défini ». Nous n'en avons pas besoin, et ce n'est assuré que si $\Omega$ est de
présentation finie.} Toute section $X$ de $\mathfrak{g}$ définit une dérivation
$\theta_X$ de $\mathcal{O}$ :
\[
\theta_X(f) = \langle \delta f, X \rangle ,
\qquad
\theta_{gX} = g\,\theta_X ,
\]
et la règle de Leibniz pour $\theta_X$ résulte de $\alpha)$. Pour $X, Y$
sections de $\mathfrak{g}$ sur un ouvert $U$, on définit $[X,Y] \in \Gamma(U,
\mathfrak{g})$ par
\[
(*)\qquad \langle \omega, [X,Y] \rangle
= -\langle \delta\omega, X \wedge Y \rangle + \theta_X \langle \omega, Y \rangle
- \theta_Y \langle \omega, X \rangle ,
\qquad \omega \in \Gamma(U, \Omega).
\]
Le second membre est $\mathcal{O}$-linéaire en $\omega$ grâce à $\alpha)$ et
$\beta)$. Il définit donc bien une section de $\mathfrak{g}$, qui est
$[X,Y]$.\footnote{La page écrit « définit une section $X \wedge Y$ » ; c'est
$[X,Y]$ qui est défini. Les deux signes devant $\theta_X$ et $\theta_Y$ dans
$(*)$ sont récrits à l'encre par-dessus d'autres, et un mot biffé illisible
suit $[X,Y]$.}

On a alors les formules suivantes.
\begin{itemize}
  \item[(**)] $\theta_{[X,Y]} = [\theta_X, \theta_Y] = \theta_X\theta_Y -
  \theta_Y\theta_X$. Plus précisément, pour toute fonction $f$,
  \[
  \theta_{[X,Y]}f - [\theta_X, \theta_Y]f = -\langle \delta\delta f, X \wedge Y\rangle ,
  \]
  de sorte que $(**)$ résulte de $\gamma)$. Réciproquement, $\gamma)$ est
  nécessaire à $(**)$ si $\Omega$ est localement libre.
  \item $[X,X] = 0$, et $[X, fY] = \theta_X(f)\,Y + f\,[X,Y]$, qui résulte de
  $\alpha)$.
  \item \emph{Jacobi.} Pour toute section $\omega$ de $\Omega$, on a l'identité
  \[
  \Big\langle \omega, \sum_{\mathrm{cycl}} [[X,Y],Z] \Big\rangle
  = \langle \delta\delta\omega, X \wedge Y \wedge Z \rangle
  - \sum_{\mathrm{cycl}} \langle \delta\delta\langle\omega, Z\rangle, X \wedge Y \rangle .
  \]
  L'identité de Jacobi vaut donc dans tout complexe de De Rham, puisque
  $\delta^{2} = 0$. Si $\Omega$ est localement libre de type fini et que
  $\gamma)$ est satisfaite, elle équivaut à la condition $\gamma')$ du
  paragraphe 2.
\end{itemize}
L'identité de Jacobi se démontre ainsi. Pour une 2-forme $\eta$, la règle de
Leibniz et $(*)$ donnent la formule de Cartan
\[
\delta\eta(X,Y,Z) = \sum_{\mathrm{cycl}} \theta_X\,\eta(Y,Z)
- \sum_{\mathrm{cycl}} \eta([X,Y],Z).
\]
On la vérifie sur les produits $a \wedge b$, sans faire appel à Jacobi. On
l'applique à $\eta = \delta\omega$, on développe $\delta\omega(Y,Z)$ par $(*)$,
et les termes en $\theta_X\langle\omega, [Y,Z]\rangle$ se détruisent par
permutation circulaire. Il reste l'identité annoncée, où $(**)$ réécrit les
termes en $[\theta_X, \theta_Y] - \theta_{[X,Y]}$.\footnote{La page mène un
calcul très corrigé, dont une partie est barrée de trois longues diagonales.
Elle aboutit à $(\mathrm{Jacobi})(\omega) = T(\delta\omega)$ avec $T =
\sum(i_X i_Y \delta i_Z - \theta_Z i_X i_Y)$. Elle récrit ensuite $T$, à l'aide
de la formule de Cartan $\theta_X = i_X\delta + \delta i_X$, sous la forme
$\sum [\,i_{X \wedge Y}, \theta_Z\,]$, et s'arrête sur : « pourquoi est-ce nul
?? ». La question reste ouverte sur la page. La réponse ci-dessus est la nôtre,
par une autre voie que la sienne. Ce qu'il fallait, c'est appliquer $\delta$ à
$\delta\omega$, c'est-à-dire utiliser $\delta^{2} = 0$ en degré 1 ; la Remarque
de la page 22 croyait cette condition automatique.}

La conclusion de la page 25 se lit alors ainsi. $\mathfrak{g}$ est une
\emph{$\mathcal{O}$-algèbre de Lie} sur $X_0$ : c'est un faisceau d'anneaux de
Lie et un $\mathcal{O}$-module. Elle est munie d'un homomorphisme d'anneaux de
Lie $\theta : \mathfrak{g} \to T_{X_0/\mathbb{Z}} =
\mathcal{D}\mathrm{er}(\mathcal{O})$, qui est $\mathcal{O}$-linéaire, et elle
satisfait $[X, fY] = f[X,Y] + \theta_X(f)\,Y$. C'est ce qu'on appelle
aujourd'hui un \emph{algébroïde de Lie} sur $X_0$, d'ancre $\theta$, ou, en
algèbre, une \emph{algèbre de Lie--Rinehart}.\footnote{La page dit
« $\mathcal{O}$-Algèbre de Lie » et « Anneau de Lie ». La notion algébrique
figure chez Rinehart (1963), après Herz et Palais ; le nom d'algébroïde de Lie
vient de Pradines (1967), en géométrie différentielle ; « algèbre de
Lie--Rinehart » est un nom plus tardif. La page ne cite aucun de ces travaux, et
nous n'impliquons pas qu'elle les connaissait.}

\pagerange{25}{26}
\subsection*{4 et 5. Le complexe de Chevalley--Eilenberg, et la correspondance (pages 25 et 26)}

Soit $\mathfrak{g}$ une $\mathcal{O}$-algèbre de Lie. On pose
$C^{p}(\mathfrak{g}) =
\mathcal{H}\mathrm{om}_{\mathcal{O}}(\Lambda^{p}_{\mathcal{O}}\mathfrak{g},
\mathcal{O})$, avec la différentielle de Chevalley--Eilenberg :
\[
(\delta c)(X_0, \ldots, X_p) = \sum_{i} (-1)^{i}\,\theta_{X_i}\,c(\ldots, \widehat{X_i}, \ldots)
+ \sum_{i<j} (-1)^{i+j}\,c([X_i, X_j], \ldots, \widehat{X_i}, \ldots, \widehat{X_j}, \ldots).
\]
En degré 1 elle redonne $(*)$.\footnote{La page renvoie à « une note Bourbaki »
et à « une formule connue », sans l'écrire. La formule, et le fait qu'elle
prolonge les conventions de signe de $(*)$, sont de nous. Pour les algèbres de
Lie--Rinehart, ce complexe est celui de Rinehart.} Lorsque $\mathfrak{g}$
provient d'un complexe de De Rham $(\Lambda^{\bullet}\Omega, \delta)$,
l'homomorphisme évident d'algèbres graduées $\Lambda^{\bullet}\Omega \to
C^{\bullet}(\mathfrak{g})$ est compatible aux différentielles. C'est un
isomorphisme si $\Omega$ est localement libre de type fini.

On obtient ainsi une correspondance bijective entre deux classes d'objets : les
$\mathcal{O}$-algèbres de Lie localement libres de type fini sur $X_0$, et les
complexes de De Rham construits sur des $\Omega$ localement libres de type
fini. On passe de l'une à l'autre par $\mathfrak{g} \mapsto
C^{\bullet}(\mathfrak{g})$ et $\Omega \mapsto
\mathcal{H}\mathrm{om}(\Omega, \mathcal{O})$. Cette correspondance renverse
les flèches : un morphisme d'algébroïdes $\mathfrak{g} \to \mathfrak{g}'$
donne un morphisme d'algèbres différentielles $\Lambda^{\bullet}\Omega' \to
\Lambda^{\bullet}\Omega$. C'est donc une \emph{anti}-équivalence.\footnote{La
page dit que les deux sortes d'objets « se correspondent exactement », sans
parler des morphismes ; la variance est notre précision. On lit aujourd'hui
cette correspondance comme celle entre algébroïdes de Lie et « champs de
vecteurs homologiques » sur le fibré décalé (Vaintrob, 1997), formulation
postérieure.}

\pagerange{26}{26}
\subsection*{6. Le théorème de Quillen (page 26)}

Revenons à une catégorie formelle $\mathcal{C}$, donnée par $A = (A_\alpha)$,
et à son complexe de De Rham $(\Lambda^{\bullet}\Omega, \delta)$.

\emph{Théorème (attribué par la page à Quillen).} Supposons $X_0$ de
caractéristique nulle, c'est-à-dire au-dessus de $\mathbb{Q}$. Alors le
foncteur $A \mapsto (\Lambda^{\bullet}\Omega_A, \delta)$ induit une
équivalence entre deux catégories :
\begin{itemize}
  \item celle des catégories formelles $I$-adiques sur $X_0$ dont le
  $\Omega$ est plat ;
  \item celle des complexes de De Rham sur $X_0$ dont le $\Omega$ est plat.
\end{itemize}
Les $A$ en question sont ceux pour lesquels le morphisme canonique
$\mathrm{Sym}_{\mathcal{O}}\Omega \to \mathrm{gr}_I A$ est un isomorphisme, ce
qu'il appelle une propriété de « lissité différentielle ».\footnote{La page ne
donne aucune référence, et ne précise pas les morphismes. Le foncteur est
covariant en $A$ : un morphisme de pro-anneaux compatible aux structures donne
un morphisme des complexes. Il est donc contravariant en $\mathcal{C}$,
puisqu'un foncteur $\mathcal{C} \to \mathcal{C}'$ induisant l'identité sur
$X_0$ donne $A' \to A$. Par le paragraphe 5, $\mathcal{C} \mapsto
\mathfrak{g}^{\mathcal{C}}$ est une équivalence \emph{covariante} vers les
algébroïdes de Lie, lorsque $\Omega$ est localement libre de type fini.
« $I$-adiques » et « lorsque $X_0$ est de caractéristique nulle » sont des
ajouts à l'encre, le second d'une lecture incertaine. La phrase sur
$\mathrm{Sym}\,\Omega \to \mathrm{Gr}(A)$ est entre crochets. Nous la lisons
comme une description des $A$ visés. Nous ne savons pas si, en
caractéristique nulle, la platitude de $\Omega$ l'entraîne en général. Pour
les groupes formels sur un corps de caractéristique nulle, la lissité est un
théorème de Cartier.}

Ceci contient, pour $X_0$ le spectre d'un corps de caractéristique nulle, le
théorème que la page attribue à Cartan. Les catégories formelles sur un point
sont alors les groupes formels, et les complexes de De Rham de dimension finie
sont les complexes de Chevalley--Eilenberg des algèbres de Lie. Le théorème dit
donc que les groupes formels de dimension finie équivalent aux algèbres de Lie
de dimension finie.\footnote{La littérature attache cette équivalence (le
« troisième théorème de Lie » formel, en caractéristique nulle) à plusieurs
noms : Cartier, Lazard, Dieudonné. Sous la forme duale des algèbres de Hopf,
c'est le théorème de Milnor et Moore (1965). Nous ne vérifions pas
l'attribution de la page. La correspondance générale entre groupoïdes formels
et algébroïdes de Lie en caractéristique nulle est aujourd'hui un chapitre
classique, dont la page ne pouvait connaître les formulations ultérieures.}

\emph{Variante à puissances divisées.} La page annonce une variante sans
hypothèse sur la caractéristique. On y remplace les voisinages infinitésimaux
par des enveloppes à puissances divisées : $A = (A_n)_{n \in \mathbb{N}}$ est
muni sur l'idéal d'augmentation d'une structure de puissances divisées, qui
passe aux quotients $A_n$. On demande que $I_n$ soit nilpotent d'ordre $n$ et
que la filtration soit celle qui provient des puissances divisées. Elle note
qu'il faudrait des relations entre cette structure et $\Delta$. L'énoncé n'est
pas formulé, et le titre « Variante à puissances divisées » qui clôt la page 30
n'est suivi de rien.\footnote{C'est le côté cristallin de la théorie : les
stratifications à puissances divisées de ses propres notes sur les cristaux
(« Crystals and the de Rham cohomology of schemes », exposé de 1966, publié
en 1968), puis les opérateurs différentiels à puissances divisées de Berthelot.
La page ne cite rien.}

\pagerange{27}{30}
\subsection*{7 à 10. Coefficients : descente et connexions (pages 27 à 30)}

\emph{Données de descente.} Soit $M$ un $\mathcal{O}$-module. On a la notion de
donnée de descente sur $M$ relativement à $\mathcal{C}$, que la page appelle
aussi « $\mathcal{C}$-stratification » ou « opération de $\mathcal{C}$ sur
$M$ ». C'est un isomorphisme $t^{*}M \simeq s^{*}M$ au-dessus de la catégorie
formelle, normalisé le long des identités et satisfaisant la condition de
cocycle. On dirait aujourd'hui un comodule sur l'algébroïde de Hopf $A$. Une
telle donnée munit $\Lambda^{\bullet}\Omega \otimes_{\mathcal{O}} M$ d'un
opérateur $\delta_M$ tel que
\[
\delta_M(a x) = \delta(a)\,x + (-1)^{|a|}\,a\,\delta_M(x)
\qquad (a \in \Lambda^{\bullet}\Omega).
\]
Il suffit pour cela de la donnée de recollement d'ordre 1. L'opérateur est
connu dès qu'on connaît sa restriction $\nabla : M \to \Omega \otimes M$. Le
fait qu'il soit de carré nul provient de la condition de descente. La page
demande si l'ordre 2 suffit pour cela, et ne répond pas.\footnote{Passage
difficile : le crochet ouvert après « La » et l'ajout interlinéaire « soit de »
se lisent mal. On en donne le sens que la transcription dégage. La page note
ici $\delta^{1}_M$ ce qu'elle note $\delta^{0}_M$ au paragraphe suivant ; nous
écrivons $\nabla$ dans les deux cas.}

\emph{Connexions relatives à un complexe de De Rham.} Plus généralement, soit
$(\Lambda^{\bullet}\Omega, \delta)$ un complexe de De Rham sur $X_0$. Une
\emph{connexion} sur $M$ est un morphisme additif $\nabla : M \to \Omega
\otimes M$ tel que $\nabla(fx) = \delta f \otimes x + f\,\nabla x$. Elle se
prolonge de façon unique en un opérateur $\delta_M$ sur
$\Lambda^{\bullet}\Omega \otimes M$ tel que
\[
\delta_M(\omega x) = \delta(\omega)\,x + (-1)^{|\omega|}\,\omega\,\delta_M(x).
\]
Le signe est nécessaire.\footnote{La page écrit $\delta_M(\omega x) =
\delta(\omega)x + \omega(\delta_M x)$ ; le signe $(-1)^{|\omega|}$ manque.}
Le carré $\delta_M^{2}$ est $\Lambda^{\bullet}\Omega$-linéaire, puisque
$\delta^{2} = 0$. Il est donc nul si et seulement s'il l'est en degré 0,
c'est-à-dire si et seulement si le composé
\[
K' : M \xrightarrow{\ \nabla\ } \Omega \otimes M \xrightarrow{\ \delta_M\ } \Lambda^{2}\Omega \otimes M
\]
est nul. Ce composé est $\mathcal{O}$-linéaire, et c'est le \emph{tenseur de
courbure}. La connexion fait opérer $\mathfrak{g}$ sur $M$ par
$\theta^{M}_X(x) = (X \otimes \mathrm{id}_M)(\nabla x)$, et l'on a
\[
a)\ \ \theta^{M}_X(fx) = \theta_X(f)\,x + f\,\theta^{M}_X(x),
\qquad
b)\ \ \theta^{M}_{fX} = f\,\theta^{M}_X ,
\qquad
c)\ \ \theta^{M}_{[X,Y]} = [\theta^{M}_X, \theta^{M}_Y] \ \text{si } K' = 0 .
\]
La condition c) vaut sous l'hypothèse d'une connexion sans
courbure.\footnote{La page biffe « intégrable » et écrit « sans courbure ».}

\emph{$\mathfrak{g}$-connexions.} Soit $\mathfrak{g}$ une
$\mathcal{O}$-algèbre de Lie. Une $\mathfrak{g}$-connexion sur $M$ est une
action $X \mapsto \theta^{M}_X$ satisfaisant a) et b). Son tenseur de courbure
est
\[
K(X,Y) = \theta^{M}_{[X,Y]} - [\theta^{M}_X, \theta^{M}_Y]
\ \in\ C^{2}\big(\mathfrak{g}, \mathcal{E}\mathrm{nd}_{\mathcal{O}}(M)\big).
\]
On vérifie que c'est bien $\mathcal{O}$-bilinéaire alterné en $(X,Y)$ et
$\mathcal{O}$-linéaire en $M$. La connexion est dite \emph{intégrable} si $K =
0$.\footnote{C'est l'opposé de la courbure usuelle $[\nabla_X, \nabla_Y] -
\nabla_{[X,Y]}$ ; la convention est celle de la page.} Si $\mathfrak{g}$
provient de $\Omega$ et que la $\mathfrak{g}$-connexion provient d'une
connexion $\nabla$, les deux courbures se déterminent l'une l'autre :
\[
K(X,Y)(u) = -\langle K'(u), X \wedge Y \rangle .
\]
C'est la même formule que celle qui relie $\theta_{[X,Y]}$ et $\delta\delta$
au paragraphe 3.\footnote{La page écrit « $K(X,Y)(u) = K'(u)(X,X)$ », où l'on
attend $(X,Y)$, et note en marge « Attention signe ! ». Le signe $-$ est celui
que donnent nos conventions ; nous l'avons vérifié.}

\emph{Le cas localement libre.} Si $\Omega$ est localement libre de type fini,
ce qui équivaut à ce que $\mathfrak{g}$ le soit, les deux notions de connexion
se correspondent biunivoquement. Il en va de même des deux tenseurs de courbure
et des deux notions de connexion intégrable. Pour une connexion intégrable, le
complexe $\Lambda^{\bullet}\Omega \otimes M$ s'identifie au complexe de
Chevalley--Eilenberg à coefficients $C^{\bullet}(\mathfrak{g}, M)$.

\pagerange{30}{30}
\subsection*{11. Cohomologie de De Rham et cohomologie de Hochschild (page 30)}

On se place sous les hypothèses du théorème de Quillen : $X_0$ de
caractéristique nulle, et $(\Lambda^{\bullet}\Omega, \delta)$ à $\Omega$ plat,
correspondant à une catégorie formelle $\mathcal{C}$ à $\Omega$ plat. Soit $M$
un $\mathcal{O}$-module. Les structures de $\mathcal{C}$-descente sur $M$
correspondent alors aux $\Omega$-connexions \emph{intégrables} sur
$M$.\footnote{La page dit « aux $\Omega$-connexions », sans l'adjectif. Il est
imposé par le paragraphe 7, où la condition de descente donne $\delta_M^{2} =
0$.} De plus, la cohomologie de De Rham de $M$, soit l'hypercohomologie
$\mathbb{H}^{\bullet}(X_0, \Lambda^{\bullet}\Omega \otimes M)$, est
canoniquement isomorphe à la cohomologie de Hochschild, ou de Čech--Alexander,
de $\mathcal{C}$ à coefficients dans $M$. Plus précisément, il y a un
quasi-isomorphisme canonique dans la catégorie dérivée
\[
\Lambda^{\bullet}\Omega \otimes M \ \simeq\ C^{\bullet}(\mathcal{C}, M).
\]
Ici $C^{n}(\mathcal{C}, M)$ est formé des fonctions sur le schéma formel des
$n$-uples de flèches composables de $\mathcal{C}$, à valeurs dans $M$, avec la
différentielle cobar.\footnote{La page n'écrit pas la définition de
$C^{\bullet}(\mathcal{C}, M)$ ; nous donnons celle qu'appellent les deux noms
qu'elle emploie. Pour $\mathcal{C}$ le complété de la diagonale d'un
$X_0$ lisse sur $\mathbb{Q}$, c'est la comparaison entre cohomologie de De Rham
et cohomologie infinitésimale par les complexes de Čech--Alexander de ses
propres notes sur les cristaux. Pour $X_0$ un point, c'est la comparaison
entre cohomologie d'un groupe formel et cohomologie de son algèbre de Lie. On
parle aujourd'hui d'isomorphisme de van Est, dont la version pour les
groupoïdes (Crainic, 2003, entre autres) est très postérieure à ces notes.}

La marge, écrite en oblique, porte : « Relier au th.\ fondamental du dévissage
des cristaux …. Peut-il y avoir ce raisonnement dans un topos convenable ?! ».
Nous n'identifions pas le théorème désigné.\footnote{La note est lue sur une vue
redressée ; sa dernière ligne est incertaine.} Suit le titre « Variante à
puissances divisées », et la page s'arrête.

Le tapuscrit du 10 septembre ne reprend pas ce paragraphe, ni rien de ces dix
pages.

\pagerange{32}{48}
\section*{Le tapuscrit du 10 septembre 1968 (pages 32 à 48 ; photocopie, pages 2 à 17)}

Le tapuscrit compte seize pages dactylographiées. Dans l'exemplaire annoté, la
page $N$ du tapuscrit est la page $31 + N$ du dossier pour $N \leq 10$, et la
page $32 + N$ pour $N \geq 11$. Dans la photocopie, c'est la page $N + 1$. Le
texte et les annotations à l'encre suivis ici sont ceux de l'exemplaire annoté.
Celui-ci porte la date « notes du 10.9.68 » en tête de sa première
page.\footnote{La même date est reproduite en tête de la photocopie, page 2.
Sur les deux exemplaires, les indices de $C_0$, $C_1$ et $X_{ij}$, frappés en
exposant, sont récrits en indice à la main, et l'on suit la correction.}

\pagerange{32}{36}
\subsection*{I. Catégories et ensembles simpliciaux (pages 32 à 36)}

\emph{Le nerf.} À toute catégorie $C$ on associe son nerf $N(C)$, l'ensemble
simplicial des suites de flèches composables. Le foncteur $N :
(\mathrm{Cat}) \to (\mathrm{sSet})$ est pleinement fidèle.\footnote{La page note
$S(C)$ le nerf et $(\mathrm{Ssimpl})$ la catégorie des ensembles
semi-simpliciaux, c'est-à-dire, dans la terminologie de l'époque, simpliciaux.
Nous réservons la lettre $S$ à la troisième partie.} Les systèmes locaux
d'ensembles sur $N(C)$ correspondent aux foncteurs sur $C$ qui envoient toute
flèche sur un isomorphisme, c'est-à-dire qui se factorisent par le groupoïde
enveloppant $C[C^{-1}]$. Leur cohomologie s'interprète de deux façons : par les
foncteurs dérivés de $\varprojlim$, ou comme cohomologie du topos $\widehat{C}$
des préfaisceaux sur $C$. Il s'agit de $H^{0}$ pour les ensembles, de $H^{1}$
pour les groupes, et de tous les $H^{i}$ pour les groupes abéliens. Pour un
système local abélien $F$, vu comme foncteur $C^{\circ} \to \mathrm{Ab}$, on a
$H^{i}(N(C), F) = H^{i}(\widehat{C}, F) = \varprojlim^{i}_{C^{\circ}} F$.

Un foncteur $u : C \to C'$ est un \emph{homotopisme} si $N(u)$ est une
équivalence faible ; c'est le terme de la page. Le critère cohomologique
d'Artin et Mazur donne une caractérisation. $u$ est un homotopisme si et
seulement si, pour tout système de coefficients local $F'$ sur $C'$, les
applications
\[
\varprojlim\nolimits^{i}_{C'^{\circ}} F' \longrightarrow \varprojlim\nolimits^{i}_{C^{\circ}} u^{*}F'
\]
sont bijectives, pour les $i$ où cela a un sens.\footnote{La page écrit au
second membre $\varprojlim^{(i)}_{C} F$, au lieu de l'image inverse de $F'$ ;
elle indexe $\varprojlim$ par $C$ plutôt que par $C^{\circ}$. La parenthèse
ouverte avant « du \emph{topos} $C$ » n'est pas refermée. « Homotopisme » se dit
aujourd'hui équivalence faible, ou équivalence d'homotopie faible, de
catégories. Le théorème A de Quillen (1973) en donnera un critère suffisant
par les catégories comma ; il est postérieur à ces notes. Le livre d'Artin et
Mazur, « Etale homotopy », paraîtra en 1969.}

\emph{Le topos.} À $C$ est associé le topos $\widehat{C}$, qui varie de façon
\emph{covariante} avec $C$. Un foncteur $u$ induit en effet un morphisme de topos
$\widehat{u} : \widehat{C} \to \widehat{C'}$, dont l'image inverse est la
restriction le long de $u$. La page ajoute que le foncteur $C \mapsto
\widehat{C}$ n'a plus rien de pleinement fidèle, « semble-t-il ».\footnote{C'est
bien le cas : le topos des préfaisceaux ne voit $C$ qu'à complétion de Karoubi
près. Deux catégories de même enveloppe karoubienne ont des topos de
préfaisceaux équivalents.} Un système de coefficients ensembliste sur $C$ est un
foncteur $C^{\circ} \to \mathrm{Ens}$ qui envoie toute flèche sur une
bijection.\footnote{La page définit les systèmes de coefficients comme les
foncteurs « transformant isomorphismes en isomorphismes », ce que fait tout
foncteur. Le sens est celui du premier alinéa : toute flèche est envoyée sur un
isomorphisme. « $C^{\circ} \to \mathrm{Ens}$ » est ajouté à l'encre.} Les
systèmes de coefficients sont exactement les objets localement constants de
$\widehat{C}$, qui se définissent intrinsèquement à partir du topos. Il en
résulte qu'être un homotopisme ne dépend que de $\widehat{u}$. Cela signifie
que $H^{i}(\widehat{C'}, F') \to H^{i}(\widehat{C}, \widehat{u}^{*}F')$ est
bijectif pour tout faisceau localement constant $F'$ sur $\widehat{C'}$.

\emph{La catégorie des simplexes.} Dans l'autre sens, on associe à un ensemble
simplicial $X$ la catégorie de ses simplexes $T(X) = \Delta/X$. Ses objets sont
les simplexes de $X$ ; c'est une catégorie fibrée sur $\Delta$, à fibres les
ensembles discrets $X_n$. Quillen prouve deux choses. D'abord, $N T(X)$ est
canoniquement isomorphe à $X$ dans $\mathrm{Ho}(\mathrm{sSet})$. Ensuite,
$T N(C)$ est canoniquement isomorphe à $C$ dans la localisée
$\mathrm{Ho}(\mathrm{Cat})$ de $(\mathrm{Cat})$ par les homotopismes. Ces
isomorphismes sont fonctoriels. Il en résulte formellement qu'un morphisme
$f$ d'ensembles simpliciaux est une équivalence faible si et seulement si
$T(f)$ est un homotopisme. $N$ et $T$ induisent donc des équivalences
quasi-inverses
\[
\mathrm{Ho}(\mathrm{Cat}) \ \simeq\ \mathrm{Ho}(\mathrm{sSet}).
\]
\footnote{C'est aujourd'hui un résultat classique : on le cite d'ordinaire
d'après Illusie (1972), qui l'attribue à Quillen. Thomason (1980) en fera une
équivalence de catégories de modèles. L'isomorphisme $N(\Delta/X) \to X$ est
réalisé par l'application « dernier sommet ».}

\emph{Une catégorie et son opposée.} Quillen construit de plus un isomorphisme
canonique et fonctoriel entre $C$ et $C^{\circ}$ dans
$\mathrm{Ho}(\mathrm{Cat})$. Il revient au même d'en construire un entre
$N(C)$ et $N(C^{\circ})$ dans $\mathrm{Ho}(\mathrm{sSet})$. C'est un
isomorphisme de la catégorie localisée, non un foncteur de $C$ dans
$C^{\circ}$.\footnote{La page ne dit pas comment Quillen le construit. Nous
n'en disons pas plus qu'elle.} Son effet sur les systèmes locaux est le suivant.
Un foncteur contravariant $F$ sur $C$, qui rend toute flèche inversible, devient
le foncteur covariant qui a les mêmes valeurs sur les objets et $F(u)^{-1}$ sur
une flèche $u$. Sur les groupoïdes fondamentaux, c'est l'isomorphisme évident
$\pi_1(C) \to \pi_1(C^{\circ}) = \pi_1(C)^{\circ}$, $g \mapsto g^{-1}$.

Il s'en déduit une interprétation faisceautique de la cohomologie d'un ensemble
simplicial $X$ à coefficients dans un système local \emph{covariant} $F$. Cette
cohomologie est classiquement définie par le complexe cosimplicial
$C^{n}(F) = \prod_{x \in X_n} F(x)$. On passe au système local contravariant
défini par $F$, on le regarde comme un faisceau sur $T(X)$, c'est-à-dire comme
un objet de $(\mathrm{sSet})_{/X}$, et on prend sa cohomologie.\footnote{La
frappe porte $\coprod_{x \in X_n}$ ; l'encre entoure le signe et demande en
marge « $\prod$ ? ». Des cochaînes sont bien un produit, et nous écrivons
$\prod$.} Pour un système de coefficients covariant non local, la page ne
connaît pas d'interprétation faisceautique de la cohomologie classique. Elle
n'en connaît pas non plus de l'homologie d'un système contravariant, ni,
inversement, de sa cohomologie faisceautique en termes classiques.

\emph{Les catégories dérivées de systèmes locaux.} Soit $u : C \to C'$ un
homotopisme. Notons $D^{b}_{\mathrm{lc}}(C)$ la sous-catégorie triangulée de
la catégorie dérivée bornée des faisceaux abéliens sur $\widehat{C}$ formée des
complexes dont les faisceaux de cohomologie sont localement constants. Quillen
montre que $u^{*} : D^{b}_{\mathrm{lc}}(C') \to D^{b}_{\mathrm{lc}}(C)$ est une
équivalence, et réciproquement. L'énoncé vaut avec un anneau de base
quelconque, supposé non nul pour la réciproque. La partie directe vaut aussi
avec un anneau de coefficients constant tordu.

Le tapuscrit conjecture ensuite (« je pense que ce résultat (facile) doit
pouvoir se généraliser ») un énoncé pour les topos. Soit $f : X \to X'$ un
morphisme de topos qui induit des isomorphismes en cohomologie pour tout
faisceau localement constant sur $X'$, le cas non commutatif inclus. Supposons
$X$ et $X'$ \emph{localement homotopiquement triviaux}. Cela veut dire que pour
tout entier $n \geq 1$, tout objet $U$ a un recouvrement $(U_i \to U)$ tel que :
a) tout système local et toute section sur $U$ deviennent constants sur les
$U_i$ ; b) pour tout groupe abélien $G$, les $H^{j}(U, G) \to H^{j}(U_i, G)$
sont nuls pour $1 \leq j \leq n$. Alors $f^{*} : D^{b}_{\mathrm{lc}}(X') \to
D^{b}_{\mathrm{lc}}(X)$ serait une équivalence. Il en irait de même avec un
système local d'anneaux sur $X'$, et $f$ induirait une équivalence entre
coefficients locaux sur $X$ et sur $X'$.\footnote{« $X$ et $X'$ », « $\geq 1$ »
sont ajoutés à l'encre, « de $X'$ » est biffé. La frappe de la fin de la page 34
a été reprise à la machine : les mots « pleinement fidèle (resp.\ » y sont
effacés, de sorte que « induit une équivalence » est la lecture la plus
probable, non la seule. Dans la phrase sur les coefficients locaux, la page
écrit $C$, $C'$ là où le contexte demande $X$, $X'$.} Une note à l'encre, en
marge, précise la portée de la condition : elle « n'est typiquement \emph{pas}
satisfaite par les schémas avec leur topologie étale (mais bien par variétés
loc.\ noeth.\ avec top.\ Zariski) ».

Le principe de démonstration proposé est le suivant. Un topos localement
homotopiquement trivial est localement connexe et localement simplement
connexe, d'où une bonne théorie des $\pi_0$ et $\pi_1$, qui sont ici discrets.
On est ramené à un cas particulier du critère d'Artin et Mazur : un
homomorphisme de groupes $G \to H$ est un isomorphisme s'il induit des
bijections sur $H^{0}$ et $H^{1}$, cas non commutatif compris. La pleine
fidélité se ramène, par la suite spectrale, aux $\mathrm{Ext}^{i}$ de
coefficients locaux. Elle résulterait de deux faits : dans un topos localement
homotopiquement trivial, les faisceaux abéliens localement constants sont
stables par $\mathcal{E}\mathrm{xt}^{i}$, et $M \mapsto M_X$ commute aux
$\mathrm{Ext}^{i}$. Pour $X$, $X'$ localement homotopiquement triviaux, trois
conditions sur $f$ seraient alors équivalentes.
\begin{itemize}
  \item[a)] $f$ est un homotopisme : il induit un isomorphisme sur les $H^{i}$
  pour tout système local sur $X'$, non nécessairement commutatif.
  \item[b)] $f^{*} : D^{b}_{\mathrm{lc}}(X') \to D^{b}_{\mathrm{lc}}(X)$ est une
  équivalence.
  \item[a')] $f$ induit une équivalence entre systèmes locaux abéliens sur $X'$
  et sur $X$, et des isomorphismes sur les $H^{i}$ correspondants.
\end{itemize}
La page ignore si l'on peut, dans a), se borner aux systèmes locaux
commutatifs, et l'encre met un point d'interrogation. La réponse est
non.\footnote{Remarque nôtre. Il existe des groupes acycliques non triviaux, par
exemple le groupe de Higman à quatre générateurs. Pour un tel groupe $G$, le
morphisme du topos $BG$ vers le point induit des isomorphismes sur la
cohomologie de tout système local abélien du point. Mais ce n'est pas un
homotopisme : $H^{1}(-, G)$ le distingue, et a') échoue. Le phénomène est celui
que la construction « plus » de Quillen, postérieure à ces notes, met en
œuvre.} L'équivalence de a) et b) motive une définition des types d'homotopie
par la catégorie $D^{b}_{\mathrm{lc}}(X)$. On la munirait de la sous-catégorie
pleine des systèmes locaux, du foncteur cohomologique à valeurs dans celle-ci et
du produit tensoriel, ce qui fait sortir de $D^{b}$ vers $D^{-}$ : « redactor
demerdetur ».\footnote{L'idée de décrire un type d'homotopie par des
coefficients reviendra chez lui, sous une tout autre forme, dans « À la
poursuite des champs » (1983). Nous n'allons pas plus loin.}

\pagerange{36}{41}
\subsection*{II. $n$-catégories, catégories $n$-uples, Gr-catégories (pages 36 à 41)}

\emph{Catégories doubles.} Le tapuscrit appelle « bi-catégorie » une catégorie
dans $(\mathrm{Cat})$. C'est la donnée d'une catégorie $C_0$, de foncteurs
source et but $p_1, p_2 : C_1 \rightrightarrows C_0$, d'un foncteur unité $C_0
\to C_1$ et d'une composition $\pi : C_1 \times_{C_0} C_1 \to C_1$, avec les
axiomes habituels. On dit aujourd'hui \emph{catégorie double}.\footnote{Ce n'est
pas une bicatégorie au sens de Bénabou (1967), c'est-à-dire une 2-catégorie
faible : c'est la catégorie double d'Ehresmann (1963). Nous employons le nom
moderne pour éviter le contresens. La page dit « un triple » pour une donnée à
cinq termes, et ne mentionne pas le foncteur unité, que nous ajoutons.} On
identifie $(\mathrm{Cat})$ à la sous-catégorie pleine de $(\mathrm{sSet})$
formée des ensembles simpliciaux qui satisfont la condition d'exactitude
« familière ». Il s'agit de transformer sommes amalgamées en produits fibrés :
$X_n \simeq X_1 \times_{X_0} \cdots \times_{X_0} X_1$, ce qu'on appelle
aujourd'hui la condition de Segal.\footnote{La page laisse la phrase en suspens
: « des objets qui …. ». Nous la complétons par la condition qu'elle énonce
elle-même une ligne plus bas. Le nom de Segal est postérieur.} De même, les
catégories doubles s'identifient aux ensembles bisimpliciaux $(X_{ij})$ qui
satisfont cette condition ligne par ligne et colonne par colonne. Les
catégories triples s'identifient de même à des objets trisimpliciaux, et
ainsi de suite.

\emph{Exemple.} Pour une catégorie $C$, posons $C_0 = C$ et $C_1 =
\mathrm{Fl}(C)$, la catégorie des flèches, avec les structures évidentes. Les
composantes simpliciales supérieures sont les catégories $\mathrm{Fl}^{n}(C)$
des suites de $n$ flèches. L'ensemble bisimplicial correspondant a pour
composantes $\mathrm{Hom}([n] \times [m], C)$. Les objets de $X_{11}$ sont les
carrés commutatifs de $C$, qu'on peut lire de deux manières comme morphismes de
flèches. L'échange des lignes et des colonnes est une symétrie de la catégorie
des catégories doubles. De même, le groupe symétrique $\mathfrak{S}_n$ opère
sur la catégorie des catégories $n$-uples.

Une \emph{2-catégorie} (stricte) est une catégorie double dont la catégorie
d'objets $C_0$ est discrète. En termes bisimpliciaux, c'est un ensemble
bisimplicial dont la première ligne est un ensemble simplicial constant. Il
paraît qu'on a une interprétation analogue des $n$-catégories en termes
d'ensembles $n$-simpliciaux. Le tapuscrit note de demander à Quillen le rapport
entre les $m$-catégories dans $(n\text{-Cat})$ et les ensembles
$(m+n)$-simpliciaux.\footnote{Paragraphe signalé par un double trait en marge.
C'est la voie que prendront, pour les $n$-catégories faibles, Tamsamani et
Simpson à la fin des années 1990, bien après ces notes.}

\emph{Groupes dans $(\mathrm{Cat})$ et catégories de Picard.} Les groupes dans
$(\mathrm{Cat})$ sont les groupes simpliciaux qui satisfont la condition de
Segal. Par la correspondance de Dold--Kan (le « tapis de Dold »), les groupes
\emph{abéliens} dans $(\mathrm{Cat})$ forment une catégorie équivalente à celle
des complexes de chaînes de longueur 1 :
\[
0 \longleftarrow C_0 \xleftarrow{\ d\ } C_1 \longleftarrow 0 .
\]
Explicitons le groupe abélien dans $(\mathrm{Cat})$ associé à un tel complexe.
On trouve la catégorie de Picard du complexe, que le tapuscrit appelle « “ma”
construction chérie ». Ses objets sont les éléments de $C_0$, et ses flèches de
$x$ vers $y$ sont les $z \in C_1$ tels que $dz = y - x$. On obtient ainsi, à
isomorphisme près, tous les groupes abéliens de $(\mathrm{Cat})$.\footnote{La
notion de catégorie de Picard est la sienne. Deligne en donnera l'exposé de
référence en 1973 (SGA 4, exposé XVIII, § 1.4). On y trouve l'équivalence entre
catégories de Picard strictement commutatives et complexes de longueur 1 à
quasi-isomorphisme près.}

Complément important : toute Gr-catégorie, c'est-à-dire toute catégorie
monoïdale où objets et flèches sont inversibles, est équivalente à un vrai
groupe dans $(\mathrm{Cat})$. Le tapuscrit disait « toute catégorie de Picard
(resp.\ et commutative) … (resp.\ $(\mathrm{Cat})$-groupe commutatif) ».
L'encre biffe les deux parenthèses, et explique en marge : une « cat.\ de
Picard commut. » ne se réalise en général que par un groupe \emph{non
commutatif} dans $(\mathrm{Cat})$.\footnote{La marge est bien fondée. Une
catégorie de Picard symétrique dont la symétrie $c_{x,x}$ n'est pas l'identité
ne peut pas être équivalente à un groupe abélien strict dans $(\mathrm{Cat})$.
Seules les catégories de Picard strictement commutatives le sont. Le terme de
Gr-catégorie, qu'on lit ici en 1968, est celui que reprendra la thèse de Hoàng
Xuân Sính (1975), préparée sous sa direction. On dit aujourd'hui aussi
2-groupes. La marge porte
en outre un premier bloc de cinq lignes, barré et resté illisible.} Une
insertion à l'encre ajoute ceci. Quand on passe aux classes d'équivalence de
catégories de Picard, cela revient à travailler dans la catégorie
\emph{dérivée} des complexes, dont les seuls invariants sont $H_0$ et $H_1$,
autrement dit $\pi_0$ et $\pi_1$.\footnote{Pour les objets, c'est exact
parce que $\mathrm{Ext}^{2}_{\mathbb{Z}} = 0$ : un complexe de longueur 1 est
quasi-isomorphe à $H_0 \oplus H_1[1]$. Les morphismes voient encore
$\mathrm{Ext}^{1}(H_0, H_1)$. « simplifie » et « Grab » sont des lectures
incertaines de l'insertion.}

Une note encadrée en tête de la page 38 se lit ainsi : « Attention : la
catégorie $\mathrm{Hom}_{\mathrm{grab}}$, qui s'identifie à la catégorie définie
par un $R\mathrm{Hom}_{\mathbb{Z}}(C_x, C'_x)$, ce point devant être élucidé !
». Les « Grab » sont les groupes abéliens ; il s'agit des homomorphismes entre
catégories de Picard associées à deux complexes $C_x$, $C'_x$. L'énoncé correct
demande une troncature. Les foncteurs additifs de $P(C_x)$ dans $P(C'_x)$
forment la catégorie de Picard associée à $\tau_{\leq 0}
R\mathrm{Hom}(C_x, C'_x)$.\footnote{La page n'appelle cette note par aucun
renvoi visible et ne définit pas $C_x$, $C'_x$. La troncature est notre
précision : $R\mathrm{Hom}$ entre complexes de longueur 1 a de la cohomologie
en trois degrés, et une catégorie de Picard n'en retient que deux. C'est sous
cette forme que l'énoncé figure chez Deligne (SGA 4 XVIII).}

On a ainsi une classification beaucoup plus restreinte que celle des
H-espaces (resp.\ des H-espaces homotopiquement commutatifs) à homotopie près.
La raison en est qu'on est plus exigeant, dans le cadre catégorique, sur
l'associativité (resp.\ la commutativité). On ne demande pas seulement
l'existence d'isomorphismes d'associativité, ce qui suffirait pour avoir un
groupe dans $\mathrm{Ho}(\mathrm{Cat}) \simeq \mathrm{Ho}(\mathrm{sSet})$. On
les suppose \emph{donnés}, avec des compatibilités, et c'est ce qui permet de
se ramener à de vrais groupes dans $(\mathrm{Cat})$. L'encre biffe la suite,
« resp.\ de vraie commutativité », pour la raison qu'on vient de dire.

\emph{Le cran suivant.} On doit trouver une meilleure approximation, d'un cran,
de la classification topologique. Il faut pour cela passer aux groupes
(resp.\ groupes abéliens) dans $(2\text{-Cat})$, ou plus généralement aux
« Gr-2-catégories » (resp.\ « Grab-2-catégories »). Deux opérations relient
les crans. Une Grab-catégorie définit une Grab-2-catégorie, comme une catégorie
définit une 2-catégorie : c'est l'analogue d'un « cosquelette tronqué ».
Inversement, une Grab-2-catégorie définit une Grab-catégorie, celle des
automorphismes de l'objet unité. Le tapuscrit y voit l'analogue du « tuage des
groupes d'homotopie en dimension $\geq 1$ ».\footnote{Les deux signes $\geq$ de
ce passage sont repassés à l'encre sur un signe frappé, et le « 1 » est
retouché. Remarque nôtre : si l'on calcule les invariants, la seconde
opération ressemble plutôt au passage à l'espace des lacets. Les $\pi_0$,
$\pi_1$ de la Gr-catégorie des automorphismes de l'unité sont les $\pi_1$,
$\pi_2$ de la Gr-2-catégorie de départ ; c'est un décalage plus qu'une
troncature.} Pour la notion « grossière », l'idée de sa lettre à Verdier était
la suivante. On doit trouver la même chose que les tronqués au cran 2 dans la
catégorie \emph{homotopique} des H-espaces, c'est-à-dire ce qu'on en déduit en
tuant les groupes d'homotopie en dimension $\geq 2$.\footnote{La lettre à
Verdier n'est pas dans le dossier.}

Le tapuscrit note de demander à Quillen si les objets-groupes et les objets
pointés connexes de la catégorie homotopique forment des catégories
équivalentes par le foncteur espace des lacets. Une encre différente répond en
marge : « non, on a un foncteur pl.\ fidèle (?) mais pas essentiellement
surjectif ». La non-surjectivité est juste. Le point d'interrogation est
justifié : le foncteur des lacets n'est même pas fidèle.\footnote{Remarque
nôtre. La classe $ab \in H^{4}(S^{2} \times S^{2}; \mathbb{Z})$, produit des
deux générateurs, définit une application non homotope à zéro $S^{2} \times
S^{2} \to K(\mathbb{Z}, 4)$. Sa bouclée est pourtant homotope à zéro, parce
que la suspension cohomologique annule les produits. « pl.\ » et
« essentiellement » sont des lectures incertaines.}

\emph{Les vrais groupes abéliens dans $(2\text{-Cat})$.} L'approche de Quillen
les identifie d'abord à des groupes abéliens dans la catégorie des catégories
doubles. Ce sont donc des groupes abéliens bisimpliciaux particuliers,
soumis à la condition de Segal sur chaque ligne et chaque colonne. Par la
version bisimpliciale de Dold--Kan, ils équivalent aux bicomplexes de groupes
abéliens nuls hors du carré $0 \leq i, j \leq 1$. La première ligne doit
correspondre à un groupe simplicial constant ; elle se réduit donc à son terme
de degré 0, et le bicomplexe devient\footnote{La frappe porte « La colonne que
la première ligne correspond » ; le sens est celui d'une condition sur la
première ligne. Les flèches du diagramme sont en partie tracées ou repassées à
l'encre.}

\begin{tikzcd}
C_{00} & 0 \arrow[l] \\
C_{10} \arrow[u] & C_{11} \arrow[l] \arrow[u]
\end{tikzcd}

Ces bicomplexes s'identifient aux complexes de chaînes de longueur 2,
$C_{11} \to C_{10} \to C_{00}$. La catégorie des groupes abéliens dans
$(2\text{-Cat})$ est donc équivalente à celle des complexes de longueur 2. La
2-catégorie associée à un tel complexe est « mon interprétation chérie ». Il
resterait à décrire, par un complexe $R\mathrm{Hom}$, la 2-Grab-catégorie des
homomorphismes entre deux telles 2-Grab-catégories strictes.

L'idée qui se dégage, et que Quillen, « on espère », va expliciter, est la
suivante. Les groupes abéliens dans $(n\text{-Cat})$ devraient correspondre
aux complexes de longueur $n$. On raisonnerait soit à isomorphisme près des deux
côtés, soit à $n$-équivalence près d'un côté et dans la catégorie dérivée de
l'autre. Cela doit pouvoir se préciser au moyen de $\mathrm{Hom}^{\bullet}$
(resp.\ $R\mathrm{Hom}^{\bullet}$).\footnote{Pour des groupes abéliens
\emph{stricts}, la catégorie dérivée est bien la bonne cible. Pour les
versions faibles et symétriques, elle ne l'est plus : ce sont des spectres
tronqués qui les classent, et non des complexes. La biffure « resp.\ de vraie
commutativité » de la page 38 enregistre déjà ce phénomène au cran 1.}

\emph{Les $K_n$ d'une catégorie monoïdale.} Soit $C$ une catégorie munie d'un
produit tensoriel associatif, et éventuellement unitaire : un bifoncteur avec
des isomorphismes d'associativité, mais sans inverses. On doit pouvoir lui
associer une Gr-catégorie universelle, comme on associe à un monoïde son
groupe symétrisé. On en tire des invariants $K_0(C)$ et $K_1(C)$. On suppose
qu'on trouvera de même, pour tout $n$, une Gr-$n$-catégorie universelle où
envoyer $C$ de façon compatible à $\otimes$. Ses invariants $\pi_n$ seraient
par définition les $K_n(C)$. Un cas intéressant est celui d'une catégorie à
produits finis, additive par exemple, où $\otimes$ est le produit. Mais pour
une catégorie abélienne, il y a une notion plus adéquate du $K_0$, qui devrait
correspondre à une meilleure notion des $K_i$, « inspirée du tapis de Roos ».
C'est sans doute en termes de « vraies » catégories triangulées qu'il faudra
poser les définitions.\footnote{Le crochet ouvert avant « Mais », tracé à
l'encre sur une parenthèse tapée, se referme au haut de la page 41. Nous
n'identifions pas le « tapis de Roos ». Quillen définira les $K_n$ supérieurs
en 1970--1972, par la construction « plus » puis par la construction $Q$
(publiée en 1973). Pour une catégorie monoïdale symétrique, ce sera par
complétion en groupe (Grayson, 1976, d'après Quillen). Schlichting (2002)
montrera que la $K$-théorie n'est pas un invariant de la catégorie
triangulée seule, ce qui donne du poids aux guillemets de « vraies ».}

Peut-être Quillen a-t-il une définition simpliciale plus directe du H-espace
qui devrait correspondre à $C$, sans passer par les $n$-catégories. En marge, à
l'encre : « oui ».\footnote{En effet, c'est par un espace que Quillen définira
les $K_n$. La marge n'est pas datée.} Ce n'est certainement pas en général le
nerf $N(C)$. Une catégorie additive est contractile, puisque le foncteur
identique est homotope au foncteur nul : il suffit d'une transformation
naturelle de l'un à l'autre. Pourtant, lorsque $C$ est une Gr-catégorie, le
nerf naïf est le bon. On trouve en effet des $\pi_0$ et $\pi_1$ qui sont aussi
ceux de $C$ au sens des catégories sans Gr-structure. Une note à l'encre au pied
de la page 41 ajoute : « A expliquer : pour des Gr-$n$-catégories, les $\pi_i$
au sens simplicial sont aussi les $\pi_i$ au sens Gr-$n$-catégories, savoir les
objets à isomorphisme près dans divers $\mathrm{Hom}(1,1)$ … ».\footnote{La
seconde moitié de cette note est barrée. Elle parlait d'une interprétation de
l'opération « lacets » en termes du décalage géométrique sur les
$n$-catégories, et le nom de l'opération est incertain.}

\pagerange{43}{48}
\subsection*{III. Point de vue « motivique » en théorie du cobordisme (pages 43 à 48)}

\emph{Le problème universel.} Soit $\mathcal{V}$ la catégorie des variétés
différentiables compactes sans bord, non nécessairement orientables. Ses
morphismes sont les classes d'homotopie d'applications continues.\footnote{La
page ne dit pas « compactes ». Nous l'ajoutons : les images directes, la
composition des correspondances et la surjectivité de Thom invoquée plus bas
l'exigent. Sur la \emph{photocopie} (page 12), et non sur l'exemplaire annoté,
une note en marge demande justement « (compactes ?) ». La généralisation de la
page 48 lève ensuite l'hypothèse.} Soit $\mathcal{E}$ une catégorie. On
s'intéresse aux couples $(F_\bullet, F^\bullet)$ formés d'un foncteur covariant
$F_\bullet : \mathcal{V} \to \mathcal{E}$ et d'un foncteur contravariant
$F^\bullet : \mathcal{V}^{\circ} \to \mathcal{E}$, qui ont les mêmes valeurs sur
les objets. On leur impose le changement de base pour tout carré cartésien
d'applications différentiables où $f$ et $g$ sont transverses :

\begin{tikzcd}
 & X' \arrow[dl, "g'"'] \arrow[dr, "f'"] & \\
X \arrow[dr, "f"'] & & Y' \arrow[dl, "g"] \\
 & Y &
\end{tikzcd}

\[
g^{\bullet} f_{\bullet} = f'_{\bullet}\, g'^{\bullet}
\qquad (f_\bullet = F_\bullet(f),\ f^\bullet = F^\bullet(f)).
\]
La question est de trouver un couple universel. Quillen prouve que la
construction suivante en donne un.\footnote{Ce cadre, deux fonctorialités liées
par le changement de base transverse, est celui que Quillen rendra célèbre en
1971 pour le cobordisme complexe (« Elementary proofs of some results of
cobordism theory using Steenrod operations »). On le retrouve dans les
théories bivariantes de Fulton et MacPherson (1981) et dans le cobordisme
algébrique de Levine et Morel (2007). Tout cela est postérieur à ces notes.}

\emph{L'anneau $B(X)$.} Pour $X$ dans $\mathcal{V}$, soit $B(X)$ l'ensemble des
classes de variétés $Z \to X$ au-dessus de $X$ pour la relation de cobordisme.
Deux variétés $Z' \to X$ et $Z'' \to X$ sont cobordantes s'il existe une
variété à bord $W$ et $W \to X$ tels que $\partial W = Z' \sqcup Z''$, les deux
applications étant induites. La somme est la réunion disjointe ; c'est une loi
de groupe où tout élément est d'ordre 2. Le produit est le produit fibré, après
qu'on a rendu les deux applications transverses par le théorème de Thom.
L'anneau est gradué par la dimension relative : $B_d(X)$ est formé des $Z \to
X$ avec $\dim Z - \dim X = d$. Pour $f : X \to Y$, l'image directe $f_\bullet :
B(X) \to B(Y)$ est la composition avec $f$, et l'image inverse $f^\bullet :
B(Y) \to B(X)$ est le produit fibré transverse. $B$ est ainsi un foncteur en
anneaux gradués contravariant, et un foncteur en groupes covariant. La formule
de changement de base entraîne la formule de projection
habituelle.\footnote{La graduation et la mention de la formule de projection
sont des ajouts à l'encre. En termes d'aujourd'hui, $B_d(X)$ est le groupe de
cobordisme non orienté $MO^{-d}(X)$ d'Atiyah (1961). Par la dualité de
Poincaré--Atiyah, c'est aussi le bordisme $\mathfrak{N}_{n+d}(X)$ de Thom, si
$X$ est de dimension $n$. La page dit « anneau de bordisme ».}

\emph{La catégorie $\mathbf{B}$.} Les constructions usuelles de la théorie des
correspondances donnent une catégorie $\mathbf{B}$. Ses objets sont ceux de
$\mathcal{V}$, et ses morphismes sont donnés par
\[
\mathrm{Hom}_{\mathbf{B}}(X, Y) = B(X \times Y),
\qquad
\beta \circ \alpha = p_{13\bullet}\big(p_{12}^{\bullet}\alpha \cdot p_{23}^{\bullet}\beta\big),
\]
où les $p_{ij}$ sont les projections de $X \times Y \times Z$.\footnote{La page
renvoie aux « constructions bien connues » ; la formule de composition est la
nôtre. Elle demande que $Y$ soit compact.} La symétrie $B(X \times Y) \simeq
B(Y \times X)$ donne un isomorphisme canonique $\mathbf{B} \simeq
\mathbf{B}^{\circ}$, qui est l'identité sur les objets. Le foncteur $F_\bullet :
\mathcal{V} \to \mathbf{B}$ est l'identité sur les objets et envoie une flèche
sur son graphe. Composé avec la symétrie, il donne $F^\bullet :
\mathcal{V}^{\circ} \to \mathbf{B}$.\footnote{La frappe écrit « $F^\bullet : C
\to \mathbf{B}$ », sans l'opposée qu'appelle un foncteur contravariant, et
l'encre n'ajoute rien.} Le couple satisfait aux conditions voulues ; la page
note que cela devrait être formalisé pour des correspondances générales.

Reste la propriété universelle. Soit $(F'_\bullet, F'^\bullet)$ un couple à
valeurs dans une catégorie $\mathbf{B}'$, et notons $X' = F'(X)$. Il faut
définir des applications $B(X \times Y) \to \mathrm{Hom}_{\mathbf{B}'}(X',
Y')$. On envoie $Z \to X \times Y$, de composantes $p$ et $q$, sur
$F'_\bullet(q) \circ F'^\bullet(p)$. Il faut voir que le résultat ne dépend que
de la classe de cobordisme de $Z$. Soit $T \to X \times Y$ un cobordisme,
avec $\partial T = Z_0 \sqcup Z_1$. En « doublant » $T$, on obtient une variété
$W$ au-dessus de $S^{1}$, transverse au-dessus de deux points $s_0$, $s_1$ dont
les fibres sont $Z_0$ et $Z_1$, et l'application vers $X \times Y$ se
prolonge à $W$. On est ramené à montrer que le composé $F'_\bullet(\alpha_i)
\circ F'^\bullet(\alpha_i)$, endomorphisme de $W'$, ne dépend pas de $i$, où
$\alpha_i : Z_i \hookrightarrow W$ est l'inclusion. Soit $p : W \to S^{1}$.
Considérons le carré cartésien d'inclusions, transverse pour $s$ valeur
régulière :

\begin{tikzcd}
W_{s} \arrow[r, "\alpha_{s}"] \arrow[d, "\alpha_{s}"'] & W \arrow[d, "{\Gamma = (p, \mathrm{id}_{W})}"] \\
W \arrow[r, "i_{s}"'] & S^{1} \times W
\end{tikzcd}

Le changement de base donne
\[
\alpha_{s\bullet}\,\alpha_s^{\bullet} = \Gamma^{\bullet}\, i_{s\bullet},
\]
et le second membre ne dépend pas de $s$, puisque la classe d'homotopie de $i_s$
n'en dépend pas. La page conclut : « Le reste est sans doute
l'AQT ».\footnote{La page écrit $X$ pour notre $W$, qui est son $\overline{T}'$.
Elle écrit la relation $\alpha_s^{*}\alpha_{s*} = \Gamma^{*} i_{s*}$. Lue comme
une composition, elle ne se compose pas, car les deux membres agissent sur des
objets différents ; nous écrivons l'ordre qui convient. Les flèches du carré,
les $\alpha$, les astérisques et l'étiquette $\Gamma = (p, \mathrm{id})$ sont
portés à l'encre dans les blancs de la frappe. Nous ne savons pas résoudre
l'abréviation « AQT ».}

\emph{La structure des $B_\bullet(X)$.} Ce sont des algèbres graduées sur
$\Lambda = B_\bullet(\mathrm{pt})$. $\Lambda$ est à degrés positifs et réduit à
$\mathbb{F}_2$ en degré 0 : c'est l'anneau de bordisme non orienté
$\mathfrak{N}_*$ de Thom. C'est une algèbre de polynômes
$\mathbb{F}_2[x_i]$, avec un générateur $x_i$ de degré $i$ pour chaque entier
$i \geq 1$ qui n'est pas de la forme $2^{\nu} - 1$.\footnote{La description des
générateurs est un ajout à l'encre, appelé par un signe en marge ; « entier
$\geq 1$ » y est d'une lecture incertaine. Elle est exacte.}

On a un homomorphisme évident d'anneaux gradués $B_\bullet(X) \to
H^{\bullet}(X)$. Il envoie $Z \to X$ sur l'image de $1 \in H^{0}(Z)$ par
l'homomorphisme de Gysin, et envoie $B_{-k}$ dans $H^{k}$ ; il est donc nul en
degrés $> 0$. Il est surjectif par Thom : toute classe provient d'une variété
compacte.\footnote{$H^{\bullet}(X)$ est la cohomologie à coefficients dans
$\mathbb{F}_2$. La surjectivité combine le théorème de représentabilité de Thom
pour l'homologie modulo 2 et la dualité de Poincaré, d'où, encore,
l'hypothèse de compacité. Sur la \emph{photocopie} (page 14), et non sur
l'exemplaire annoté, une note précise « NB $H^{*}(X)$ ($= H^{*}(X,
\mathbb{F}_2)$) ».} Une note à l'encre ajoute : « en fait $H^{\bullet}(X)
\simeq B_\bullet(X)/\Lambda^{+}B_\bullet(X)$ », et c'est exact.

On prouve même qu'il existe un relèvement fonctoriel $s : H^{\bullet}(X) \to
B_\bullet(X)$, tel que l'homomorphisme induit
\[
H^{\bullet}(X) \otimes_{\mathbb{F}_2} \Lambda \xrightarrow{\ \sim\ } B_\bullet(X)
\]
soit un isomorphisme de $\Lambda$-modules. Le tapuscrit ajoute que ce sera
vrai pour tout relèvement pris individuellement, et c'est exact pour tout
relèvement homogène.\footnote{Justification (nôtre). $B_\bullet(X)$ est un
$\Lambda$-module libre (Conner et Floyd, 1964), de quotient $H^{\bullet}(X)$
modulo $\Lambda^{+}$. Par le lemme de Nakayama gradué, les images d'une base
homogène engendrent, et un épimorphisme entre modules libres de même rang
gradué, de dimension finie en chaque degré, est un isomorphisme.} Le
tapuscrit disait le relèvement « compatible avec multiplication ». L'encre
écrit au-dessus « il est douteux qu'il soit ». Elle transforme aussi la seconde
condition en question : « Quillen ignore si » le relèvement est compatible avec
les homomorphismes de Gysin.\footnote{Les lettres a) et b) et le « et » sont
biffés à la main, et la seconde clause est mise entre crochets à l'encre. Ces
deux questions, multiplicativité et compatibilité aux images directes, sont
exactement celles que tranchera la note de Quillen de 1969, « On the formal
group laws of unoriented and complex cobordism theory ». La loi de groupe
formel du cobordisme non orienté y est universelle parmi les lois $F$ sur les
$\mathbb{F}_2$-algèbres telles que $F(x,x) = 0$. Toute telle loi est isomorphe
à la loi additive, mais n'est pas égale à elle, et c'est cet écart qui
gouverne la compatibilité d'un relèvement multiplicatif aux images de Gysin.
Une formule de Riemann--Roch en tient lieu. Ces notes du 10 septembre 1968 ne
pouvaient pas le savoir, et nous ne leur prêtons rien de ce qui précède.}

\emph{Modules et enveloppe de Karoubi.} Le foncteur $X \mapsto B_\bullet(X)$
est pleinement fidèle, de $\mathbf{B}$ dans la catégorie des
$\Lambda$-modules gradués libres de type fini. Les modules obtenus ont leurs
générateurs en degrés $\leq 0$. Pour $X$ connexe, il y a un seul générateur de
degré 0, et ils satisfont la dualité de Poincaré.\footnote{La pleine fidélité
résulte de l'isomorphisme précédent, de la formule de Künneth et de la dualité
de Poincaré ; la page ne la détaille pas.} Passons à l'enveloppe
pseudo-abélienne $\widetilde{\mathbf{B}}$ de $\mathbf{B}$, c'est-à-dire à son
enveloppe de Karoubi, puisque $\mathbf{B}$ est déjà additive. On obtient tous
les $\Lambda$-modules gradués libres de type fini à générateurs en degrés $\leq
0$. C'est une catégorie additive et karoubienne, munie d'un $\otimes$ induit par
le produit des variétés. Elle a toutes les vertus, sauf d'être abélienne et
d'avoir des $\mathrm{Hom}$ internes.\footnote{La page dit « catégorie
quasi-abélienne », au sens de pseudo-abélienne. Ce n'est pas la catégorie
quasi-abélienne au sens moderne (Schneiders, 1999), notion toute différente.
Le tilde de $\widetilde{\mathbf{B}}$ est à l'encre. Sur la \emph{photocopie}
(pages 14 et 15), et non sur l'exemplaire annoté, deux notes le disent en
toutes lettres : « i.e.\ l'enveloppe “karoubienne” » et « quasi-abélienne $=$
karoubienne $+$ additive ».}

\emph{L'objet $\Lambda(-1)$.} Pour obtenir les $\mathrm{Hom}$ internes, on peut
paraphraser la construction des motifs. Soit $p : S^{1} \to e$ l'application du
cercle sur un point. Le morphisme $p^{\bullet} : e \to S^{1}$ de $\mathbf{B}$
est scindé par le choix d'un point de $S^{1}$. On pose
\[
\Lambda(-1) = \mathrm{coker}\big(p^{\bullet} : e \to S^{1}\big)
\ \in\ \widetilde{\mathbf{B}},
\qquad
B_\bullet(S^{1}) = \Lambda \oplus \Lambda(-1).
\]
C'est un module libre sur un générateur de degré $-1$, donc un objet de poids
1, le poids étant le degré des générateurs changé de signe.\footnote{La frappe
écrit $S$ pour le cercle, et « $\Lambda(-1)$ » est ajouté à l'encre. Le
scindage, qui fait exister le conoyau dans l'enveloppe karoubienne, est notre
précision. C'est l'analogue de la décomposition du motif de la droite
projective en l'unité et le motif de Lefschetz.} Le foncteur $M \mapsto M(-1) =
M \otimes \Lambda(-1)$ est pleinement fidèle. Une construction formelle
standard permet donc de rendre $\Lambda(-1)$ inversible. Dans la catégorie
obtenue, les $\mathrm{Hom}$ internes existent et valent $\check{X} \otimes Y$.
On y trouve \emph{tous} les $\Lambda$-modules gradués libres de type fini.

Cette catégorie n'est pas abélienne, puisque $\Lambda$ n'est pas un corps. On
s'attend à ce que la catégorie abélienne engendrée soit celle des
$\Lambda$-modules gradués de présentation finie. Cette catégorie est bien
abélienne, quoique $\Lambda$ ne soit pas noethérien. $\Lambda$ est en effet
limite inductive d'anneaux de polynômes en un nombre fini de variables, à
morphismes de transition plats ; il est donc cohérent.\footnote{Le tapuscrit
demandait si $\Lambda$ était noethérien, et si le nombre de générateurs était
fini. La question est biffée, et l'encre répond par la phrase que nous
rendons. « de type » est corrigé en « de présentation » à l'encre. Le « lim »
de l'encre est souligné d'un trait qui peut être la flèche d'une limite
inductive. Le terme d'anneau cohérent est le nôtre.}

\emph{L'analogie avec les motifs.} Cette construction ressemble à une théorie
des motifs où l'on n'aurait pas passé à l'équivalence homologique des cycles,
mais gardé l'équivalence rationnelle, bien trop fine à beaucoup
d'égards.\footnote{On parle aujourd'hui de motifs de Chow, dans le cadre exposé
par Manin en 1968.} On pourrait poser le même problème universel en prenant
l'équivalence homologique des morphismes, via le graphe, au lieu de
l'équivalence homotopique. Un passage biffé en tirait une conclusion
décevante. La catégorie solution aurait pour flèches toutes les correspondances
cohomologiques, et donnerait, après inversion de $\mathbb{F}_2(-1)$, tous les
espaces vectoriels gradués de dimension finie sur $\mathbb{F}_2 $ : « un peu
trop trivial ».\footnote{La dernière ligne de la page 46 et les sept premières
de la page 47 sont barrées.}

\emph{La catégorie $\mathbf{S}$ et le groupe de Steenrod.} Pour trouver une
catégorie universelle, on associe à $X$ l'ensemble $S(X) \subset
H^{\bullet}(X)$ des classes de sous-variétés de $X$, modulo équivalence
homologique. On prouve, par un théorème de Thom dont il va être question, que
c'est un sous-anneau. Par le même théorème, les $S(X \times Y)$ sont stables par
composition des correspondances cohomologiques. Cela ne semble pas trivial,
puisque les $S(X)$ ne sont pas stables par Gysin. D'où une catégorie additive
$\mathbf{S}$, dont les objets sont les variétés et les flèches les $S(X \times
Y)$. « Il semble qu'elle doive jouer le rôle universel que j'ai envisagé plus
haut. »

Pour en donner une interprétation algébrique, considérons l'algèbre de
Steenrod $\mathcal{A}$ des opérations cohomologiques stables modulo 2. Elle est
graduée, et munie d'une diagonale associative et unitaire, anticommutative,
c'est-à-dire commutative (on est en caractéristique 2). Son dual gradué
$\mathcal{A}_*$ est une algèbre commutative, avec une diagonale associative et
unitaire, et réduite à $\mathbb{F}_2$ en degré 0. C'est donc l'anneau d'un
schéma en groupes affine $G$ sur $\mathbb{F}_2$. La page l'identifie au foncteur
des $\otimes$-automorphismes additifs de $H^{\bullet}$ sur $\mathcal{V}$ qui
opèrent trivialement sur $H^{1}(S^{1})$.\footnote{« i.e.\ commutative (car.\ 2
!) », « graduée anticommutative donc » et « additifs » sont frappés en
interligne. La définition tapée de l'algèbre des opérations est barrée. La
restriction « qui opèrent trivialement sur $H^{1}(S^{1})$ » est à l'encre, et
l'exposant 1 est incertain. C'est elle qui écarte les automorphismes de
graduation. Aujourd'hui, $\mathcal{A}_* = \mathbb{F}_2[\xi_1, \xi_2, \ldots]$
(Milnor, 1958), et $G(R)$ est le groupe, pour la composition, des séries $x +
\sum a_i x^{2^{i}}$ : les automorphismes stricts du groupe formel additif.}

Le tapuscrit énonce alors, comme « le théorème de Thom », qu'un $x \in
H^{\bullet}(X)$ est dans $S(X)$ si et seulement s'il est fixe par $G$,
c'est-à-dire si $\mathcal{A}^{+} x = 0$. Il demande aussitôt : « Pourquoi cette
condition est-elle déjà nécessaire ? ». Nous rapportons cet énoncé sous réserve.
Lu avec l'action usuelle des carrés de Steenrod sur la cohomologie, il est en
défaut.\footnote{Remarque nôtre. Une droite $\mathbb{RP}^{1} \subset
\mathbb{RP}^{2}$ est une sous-variété, de classe le générateur $a \in
H^{1}(\mathbb{RP}^{2})$, et $\mathrm{Sq}^{1}a = a^{2} \neq 0$. La nécessité
échoue donc, comme la page le soupçonne. L'énoncé doit viser une autre action,
par exemple sur l'homologie ou tordue par les classes de Wu, que la page ne
permet pas de restituer. Ce qui en dépend ci-dessous est rapporté sous la
même réserve.} Sous cette réserve, le foncteur $H^{\bullet}$, de
$\mathbf{S}$ vers les représentations de $G$ dans des espaces vectoriels sur
$\mathbb{F}_2$, serait pleinement fidèle. Son enveloppe pseudo-abélienne
serait une sous-catégorie pleine des représentations de $G$. Il est douteux
qu'elle soit abélienne, et plus encore qu'elle contienne toutes les
représentations de dimension finie.

$G$ est prolisse, et, « d'après Quillen », pro-unipotent. Une note à l'encre en
marge répond « oui, car son anneau affine est une algèbre de polynômes à une
infinité d'indéterminées, réunion d'algèbres de polynômes » en un nombre fini
d'indéterminées, stables par $\Delta$.\footnote{La fin de la note est en partie
illisible. Le « car » justifie la prolissité, par les sous-algèbres
$\mathbb{F}_2[\xi_1, \ldots, \xi_n]$ stables par la diagonale. La
pro-unipotence vient de ce que l'algèbre de Hopf est graduée connexe, réduite à
$\mathbb{F}_2$ en degré 0.} Ici $G$ joue le rôle d'un \emph{groupe de Galois
motivique}. Contrairement à la situation familière, il est fortement non
réductif, et son corps de définition est $\mathbb{F}_2$, non
$\mathbb{Q}$.\footnote{L'expression « groupe de Galois motivique » est sur la
page, en 1968. Le formalisme tannakien qui lui donne son sens précis sera
publié par Saavedra Rivano en 1972. Nous ne tirons de cette rencontre aucune
conclusion de priorité.}

\emph{Orientations.} Quillen généralise la construction de $\mathbf{B}$ de deux
façons. D'abord aux variétés non nécessairement compactes, $f_\bullet$ n'étant
alors défini que pour $f$ propre. Ensuite à une condition d'orientation stable
sur les \emph{morphismes}, donnée par un morphisme $H \to BO(\infty)$, où $H$
est un H-espace, dont la page se demande s'il est commutatif, associatif et
unitaire. Une orientation de $f : X \to Y$ est un relèvement homotopique à $H$
de l'application $X \to BO$ qui classe $T_X - f^{\bullet}T_Y$. Les morphismes
$\mathrm{Hom}_{\mathbf{B}^{H}}(X, Y)$ sont les classes de diagrammes $X
\leftarrow Z \to Y$ de morphismes orientés, le second propre, modulo les bords.
Comme ensemble, cela se décrit par la variété $X \times Y$ munie de la famille
des fermés propres sur $Y$.\footnote{La page note $H(X,Y)$ ces ensembles de
morphismes, avec la même lettre que l'H-espace ; nous changeons la notation.
Pour $H = BO$ on retrouve $\mathbf{B}$. Pour $H = BU$, on obtient le
cobordisme complexe $MU$, celui où Quillen trouvera en 1969 la loi de groupe
formel universelle. Ce sont les $(B,f)$-structures de Lashof (1963).} Une note à
l'encre au pied de la page 48 ajoute : « Les $B^{H}_m(X)$ ($X$ compact de dim.\
$n$) peuvent se décrire homotopiquement à la Thom, comme les groupes
d'homotopie stables d'espaces $\mathrm{Hom}(X, MH(N - m))$ ». C'est la
construction de Pontryagin--Thom, $MH$ désignant l'espace de Thom associé à
$H$.\footnote{L'indexation est celle de la page, qui ne l'explique pas. La
note se termine par « - - - - - ».}

\pagerange{2}{17}
\section*{La photocopie et ses notes propres (pages 2 à 17)}

La photocopie des pages 2 à 17 reproduit en noir l'exemplaire annoté, encre
comprise. Elle a donc le même texte et la même couche manuscrite que les pages
32 à 48. Elle porte en outre cinq notes qui ne sont pas sur l'exemplaire annoté.
On les a rencontrées en note dans la troisième partie :
\begin{itemize}
  \item page 12, en regard des « variétés différentiables » : « (compactes ?)
  », qui est l'hypothèse que nous avons dû ajouter ;
  \item page 14 : « NB $H^{*}(X)$ ($= H^{*}(X, \mathbb{F}_2)$) », qui fixe
  les coefficients ;
  \item page 14, en regard de la « catégorie quasi-abélienne » : « i.e.\
  l'enveloppe “karoubienne” » ;
  \item page 15 : « quasi-abélienne $=$ karoubienne $+$ additive » ;
  \item page 17, au pied de la dernière page : « cf exposé Karoubi à Bourbaki !
  ».
\end{itemize}
Les trois premières précisent des hypothèses que le texte laisse implicites.
Les deux suivantes traduisent le vocabulaire du tapuscrit dans celui qui
s'imposera. La dernière pointe, selon toute vraisemblance, vers l'exposé de
Karoubi au Séminaire Bourbaki de 1971-1972, « Cobordisme et groupes formels
(d'après D. Quillen et T. tom Dieck) ». Il rend compte des résultats de Quillen
sur le cobordisme et les lois de groupes formels, ceux-là mêmes qui répondent
aux questions de la page 45. Si c'est bien cet exposé, les notes propres à la
photocopie lui sont postérieures. C'est une
inférence, non une date lue ; la datation du dossier reste celle de
l'inventaire.

\pagerange{50}{68}
\section*{Les pièces de Quillen (pages 19, 50 à 68)}

Trois pièces sont de Quillen lui-même. Elles ne sont pas transcrites, faute de
porter autre chose que ce que Quillen a publié ou diffusé.
\begin{itemize}
  \item \emph{Page 19} : la liste dactylographiée « Publications of Daniel
  Quillen », onze titres, de 1966 à « to appear ». Elle ne porte aucune
  annotation.
  \item \emph{Pages 50 à 55} : le tiré à part « On the Representation of
  Hermitian Forms as Sums of Squares », Inventiones math. 5 (1968), 237--242.
  Il ne porte aucune marque de main.
  \item \emph{Pages 57 à 68} : le tapuscrit anglais « On a conjecture of
  Adams », reproduit à l'alcool, ses pages 1 à 12, avec l'écriture de Quillen
  reproduite par la duplication. Sa seule intervention étrangère est à la page
  58, d'une encre plus foncée. La phrase porte sur l'application canonique $F :
  E \to E^{(p)}$ : elle est « of (inseparable) degree $p$ on each fiber ». On
  a mis $\dim E$ en exposant de $p$, et le degré devient $p^{\dim E}$, ce qui
  est juste : le Frobenius relatif d'un fibré vectoriel de rang $r$ est de
  degré $p^{r}$.\footnote{La main n'est pas identifiée. Les pages 61 à 68 ne
  portent rien de tel, et ne sont pas transcrites.}
\end{itemize}

\end{document}
